\chapter{Revisão Estruturada da Literatura}
\label{cap:revisao-literatura}

Esta seção apresenta o panorama do estado da arte sobre o uso de testes de mutação para avaliar a eficácia de suítes de testes automatizados, com ênfase no contexto de aplicações web. A investigação foi conduzida por meio de uma Revisão Sistemática da Literatura (RSL), estruturada para mapear as principais abordagens, métricas comparativas, ferramentas utilizadas e lacunas científicas existentes na área.

\section{Planejamento e Protocolo}
\label{sec:planejamento-protocolo}

O planejamento de uma revisão sistemática define, antes da busca, o objetivo, as questões de pesquisa, a estratégia de busca e os critérios de seleção e de extração dos dados, de modo que o processo possa ser reproduzido \cite{kitchenham2007}. Esta seção apresenta o protocolo seguido nesta revisão, cujos resultados são discutidos na Seção~\ref{sec:resultados}.

\section{Estruturação da String de Busca}
\label{subsec:string-busca}

Para estruturar a delimitação da busca na base digital, utilizou-se o protocolo PICO, amplamente empregado na área da medicina e adaptado para investigações em Engenharia de Software \cite{kitchenham2007}. O acrônimo PICO refere-se a \textit{Population} (população), \textit{Intervention} (intervenção), \textit{Comparison} (comparação) e \textit{Outcome} (resultado), podendo ser conceitualmente representado por um diagrama de Venn, conforme ilustrado na Figura~\ref{fig:pico-venn} \cite{pai2004}. Sob essa perspectiva conceitual, os estudos mais relevantes para a pesquisa são aqueles localizados precisamente na intersecção lógica entre os conjuntos definidos pelos elementos do protocolo.

Considerando o escopo investigado nesta pesquisa, a população compreende o domínio de software, servindo de base para a abrangência dos estudos mapeados. A intervenção refere-se à técnica principal avaliada, sendo contemplados os testes de mutação e a injeção sintética de falhas. O resultado representa a saída esperada e relatada nos estudos selecionados, focando em mensurar a eficácia da suíte de testes automatizados e sua capacidade de detecção de defeitos. Quanto à comparação, embora o modelo PICO original preveja a análise comparativa com tratamentos alternativos, optou-se por focar a investigação na análise da eficácia da técnica em relação às métricas de teste, dispensando a utilização desta camada específica na formulação da string, conforme flexibilizações metodológicas comuns na área.

Desta forma, os termos de busca foram gerados a partir do mapeamento dos eixos do PICO e da seleção das expressões utilizadas na literatura para se referir a cada dimensão investigada. A estrutura consolidada contendo os elementos do PICO, as palavras-chave primárias e os seus respectivos sinônimos e termos relacionados é apresentada na Tabela~\ref{tab:termos-string-busca}.

\begin{table}[htb]
    \centering
    \caption{Estrutura dos Termos da String de Busca}
    \label{tab:termos-string-busca}
    \begin{tabularx}{0.95\textwidth}{|l|p{3.8cm}|X|}
        \hline
        \textbf{Elementos} & \textbf{Palavras Chaves} & \textbf{Sinônimos e Termos Relacionados} \\ \hline
        População & software & software systems, software applications, computer software, software products, software systems development, software engineering, web app, web system \\ \hline
        Intervenção & testes de mutação & mutation testing, mutation analysis, mutation-based testing, mutation testing techniques, mutation operators, artificial faults, artificial changes, seeded faults, mutant generation, software mutation \\ \hline
        Resultados & eficácia da suíte de testes automatizados & test suite effectiveness, test suite adequacy, automated test effectiveness, mutation score, mutant detection rate, fault detection, defect detection, test effectiveness, test adequacy, test coverage \\ \hline
    \end{tabularx}
    \fonte{Autor(2026).}
\end{table}

Para finalizar a construção da expressão de busca, adicionou-se o operador lógico \texttt{AND} para separar cada camada do protocolo e, dentro de cada uma delas, o operador lógico \texttt{OR} para agrupar as palavras-chave e seus sinônimos. Acrescentou-se também a sintaxe de aspas duplas para o agrupamento exato de termos compostos. Isto resultou na seguinte estruturação da string finalizada:


\vspace{0.3cm}
\noindent\textit{( "complex systems" OR "software industry" OR "software systems" OR "software applications" OR "computer software" OR "software products" OR "software systems development" OR "software engineering" OR "web app" OR "web system" ) AND ( "mutation testing" OR "mutation analysis" OR "mutation-based testing" OR "mutation testing techniques" OR "mutation operators" OR "artificial faults" OR "artificial changes" OR "seeded faults" OR "mutant generation" OR "software mutation" ) AND ( "test suite effectiveness" OR "test suite adequacy" OR "automated test effectiveness" OR "mutation score" OR "mutant detection rate" OR "fault detection" OR "defect detection" OR "test effectiveness" OR "test adequacy" OR "test coverage" OR "test quality" )}
\vspace{0.3cm}

\begin{figure}[H]
    \centering
    \caption{Representação Visual do Protocolo PICO}
    \label{fig:pico-venn}
    \fbox{\includegraphics[width=0.8\textwidth]{figuras/PICO.png}}
    \vspace{0.2cm}
   \fonte{Adaptado de \citeonline{pai2004}.}
\end{figure}

\section{Realização da Busca e Seleção dos Artigos}
\label{sec:realizacao-busca}

Com a string de busca devidamente estruturada, procedeu-se à execução da busca automatizada na base de dados digital Scopus no dia 17 de setembro de 2026, retornando um total de 238 artigos brutos. Em seguida, aplicou-se um filtro temporal para selecionar apenas publicações do período de 2010 até o presente momento, o que reduziu a amostragem para 179 artigos.

Esses registros foram importados para a ferramenta Evidentia, ambiente utilizado para o gerenciamento do protocolo da RSL. Nessa etapa de importação, identificou-se que 4 arquivos encontravam-se corrompidos, restando um corpus inicial de 175 artigos válidos para análise.

A partir desse conjunto, estabeleceram-se os critérios formais de inclusão e exclusão para orientar triagem dos estudos, conforme sintetizado na Tabela~\ref{tab:criterios-inclusao-exclusao}.

\begin{table}[htb]
    \centering
    \caption{Critérios de Inclusão e Exclusão}
    \label{tab:criterios-inclusao-exclusao}
    \begin{tabularx}{0.95\textwidth}{|p{0.46\textwidth}|p{0.46\textwidth}|}
        \hline
        \textbf{Critérios de Inclusão} & \textbf{Critérios de Exclusão} \\ \hline
        1. Estudos que envolvam suítes de testes automatizados utilizados para detectar mutantes. & 1. Estudos que não abordem testes de mutação. \\ \hline
        2. Estudos aplicados a aplicações web, sistemas web ou software desenvolvido para a web. & 2. Estudos que não apresentem resultados relacionados à detecção de mutantes, eficácia ou adequação dos testes. \\ \hline
        3. Artigos publicados em periódicos ou conferências científicas. & 3. Trabalhos que sejam teses, dissertações, trabalhos de conclusão de curso, livros, capítulos de livros, resumos, pôsteres ou tutoriais. \\ \hline
        4. Estudos publicados nos idiomas inglês ou português. & 4. Estudos duplicados encontrados em diferentes bases de dados. \\ \hline
        5. Estudos cujo texto completo esteja disponível para análise. & \multicolumn{1}{c|}{--} \\ \hline
    \end{tabularx}
    \fonte{Autor (2026).}
\end{table}

Em um primeiro momento, aplicaram-se os critérios de inclusão e exclusão mediante a leitura preliminar do título, resumo e palavras-chave de cada um dos 175 artigos. Os estudos que atenderam aos critérios nessa triagem inicial totalizaram 42 artigos, os quais foram selecionados para a etapa de leitura aprofundada na íntegra.

Posteriormente, realizou-se a leitura completa desses 42 artigos. Durante essa análise detalhada, identificou-se que 5 estudos não apresentavam compatibilidade direta ao escopo da pesquisa após a leitura integral, sendo descartados. Dessa forma, restou um corpus final de 37 artigos aprovados.

Os dados relevantes desses 37 artigos foram extraídos com o objetivo de responder às perguntas de pesquisa secundárias da investigação, apresentadas na Tabela~\ref{tab:perguntas-secundarias}.

\begin{table}[htb]
    \centering
    \caption{Perguntas Secundárias}
    \label{tab:perguntas-secundarias}
    \begin{tabularx}{0.95\textwidth}{|l|X|}
        \hline
        \textbf{Item} & \textbf{Perguntas Secundárias} \\ \hline
        PS1 & Quais métricas são mais utilizadas na literatura para quantificar a eficácia das suítes de testes em aplicações web? \\ \hline
        PS2 & Quais operadores de mutação e tipos de alterações artificiais no código-fonte são mais aplicados no contexto específico de aplicações web? \\ \hline
        PS3 & Quais ferramentas automatizadas de teste de mutação para tecnologias web são reportadas nos estudos primários e qual o seu nível de maturidade? \\ \hline
        PS4 & Quais são as principais limitações e desafios práticos identificados na aplicação de testes de mutação em aplicações web? \\ \hline
        PS5 & Quais estratégias são propostas na literatura para reduzir os custos e mitigar as limitações da aplicação de testes de mutação em aplicações web? \\ \hline
        PS6 & Quais fatores são associados à sobrevivência de mutantes nas suítes de testes avaliadas? \\ \hline
        PS7 & Qual é a distribuição, por ano, das publicações sobre testes de mutação aplicados à avaliação de suítes de testes automatizados em aplicações web? \\ \hline
        PS8 & Em quais revistas e conferências são publicados estudos sobre o tema? \\ \hline
        PS9 & Quais métodos de validação empírica são utilizados nos estudos? \\ \hline
        PS10 & Quais são os principais autores e grupos de pesquisa que investigam o tema? \\ \hline
    \end{tabularx}
    \fonte{Autor (2026).}
\end{table}

Para padronizar a coleta das evidências durante a análise, utilizou-se um formulário de extração de dados estruturado, cujos campos e diretrizes constam na Tabela~\ref{tab:formulario-extracao}.

\begin{table}[htb]
    \centering
    \caption{Formulário de Extração dos Dados}
    \label{tab:formulario-extracao}
    \begin{tabularx}{0.95\textwidth}{|l|X|}
        \hline
        \textbf{Item} & \textbf{Descrição} \\ \hline
        1 & Título \\ \hline
        2 & Ano de Publicação \\ \hline
        3 & Fonte de Publicação \\ \hline
        4 & Autores \\ \hline
        5 & Palavras-chaves \\ \hline
        6 & Métricas utilizadas para avaliar a eficácia \\ \hline
        7 & Operadores mutantes utilizados \\ \hline
        8 & Ferramentas utilizadas/reportadas \\ \hline
        9 & Limitações e desafios \\ \hline
        10 & Estratégias para redução de limitações \\ \hline
        11 & Fatores associados à sobrevivência de mutantes \\ \hline
        12 & Método de validação empírica \\ \hline
    \end{tabularx}
    \fonte{Autor (2026).}
\end{table}

\section{Resultados}
\label{sec:resultados}

Após a conclusão da leitura dos artigos aprovados, os dados foram extraídos utilizando a ferramenta Evidentia, ambiente adotado para conduzir o protocolo da revisão. Depois de finalizada a extração dos dados, as informações coletadas foram importadas para uma planilha eletrônica \href{https://docs.google.com/spreadsheets/d/1vhQH-s1NbXxbqH6Kie3z7ZXoASvZnZw09Wph0cionCQ/edit?usp=sharing}{planilha eletrônica}, permitindo consolidar os registros, garantir a rastreabilidade e facilitar o cruzamento dos dados.

Com a planilha devidamente estruturada, procedeu-se à análise quantitativa e qualitativa do corpus de 37 artigos para responder às questões de pesquisa secundárias estabelecidas no planejamento desta investigação. Os principais resultados obtidos são apresentados e discutidos nas subseções a seguir.

\subsection[Quais métricas são mais utilizadas na literatura?]{Quais métricas são mais utilizadas na literatura para quantificar 
a eficácia das suítes de testes em aplicações web?}
\label{subsec:metricas-utilizadas}

O \textit{mutation score} é a métrica mais utilizada: 29 dos 37 estudos (81\%) o adotam ou usam uma variante diretamente equivalente, como o \textit{kill score}, a cobertura de mutação ou a contagem de mutantes mortos. Em seguida aparecem as contagens e taxas de mutantes mortos, vivos, equivalentes e teimosos (12 estudos, 33\%) e as métricas de cobertura de código (9 estudos, 25\%). Com menor frequência surgem métricas de qualidade da seleção de mutantes, de custo e de validade frente a falhas reais. Apenas dois estudos não usam o \textit{mutation score} nem derivados dele: o CBUA usa somente cobertura de código \cite{zhang2022} e o DeepMutation reporta apenas métricas da própria ferramenta \cite{tufano2020}.

Uma ressalva importante: apenas um dos 37 estudos avalia suítes de testes de aplicações web \cite{yandrapally2021}. As frequências da Tabela~\ref{tab:metricas-utilizadas} descrevem a literatura de teste de mutação em geral, majoritariamente sobre programas Java e C. 

\begingroup
\footnotesize
\setlength{\tabcolsep}{4pt}
\begin{longtable}{|p{0.28\textwidth}|>{\centering\arraybackslash}p{0.08\textwidth}|>{\centering\arraybackslash}p{0.07\textwidth}|p{0.51\textwidth}|}
    \caption{Métricas utilizadas para quantificar a eficácia das suítes de testes}
    \label{tab:metricas-utilizadas} \\
    
    % --- Cabeçalho da PRIMEIRA página ---
    \hline
    \textbf{Categoria de métrica} & \textbf{Estudos} & \textbf{\%} & \textbf{Referências} \\
    \hline
    \endfirsthead

    % --- Cabeçalho das PÁGINAS SEGUINTES ---
    \multicolumn{4}{c}{\footnotesize\textit{Continuação da Tabela~\ref{tab:metricas-eficacia}}} \\
    \hline
    \textbf{Categoria de métrica} & \textbf{Estudos} & \textbf{\%} & \textbf{Referências} \\
    \hline
    \endhead

    % --- Rodapé intermediário ---
    \hline
    \endfoot

    % --- Rodapé FINAL ---
    \hline
    \multicolumn{4}{c}{\footnotesize Fonte: Autor (2026).} \\
    \endlastfoot

    % --- Conteúdo da tabela ---
    \textit{Mutation score} e variantes diretas (\textit{kill score}, cobertura de mutação, contagem de mutantes mortos) & 29 & 81\% & Ver o texto desta seção \\
    \hline
    Contagens e taxas de mutantes mortos, vivos, equivalentes e teimosos & 12 & 33\% & \citeonline{just2012}; \citeonline{esterobotaro2010}; \citeonline{belli2011}; \citeonline{guerino2024}; \citeonline{lin2018}; \citeonline{petrovic2021}; \citeonline{arasteh2023}; \citeonline{yandrapally2021}; \citeonline{kintis2010}; \citeonline{yao2014}; \citeonline{assylbekov2013}; \citeonline{arasteh2022} \\
    \hline
    Cobertura de código (linha, instrução, ramo e método) & 9 & 25\% & \citeonline{nica2011}; \citeonline{assylbekov2013}; \citeonline{yao2014}; \citeonline{just2014}; \citeonline{ahmed2016}; \citeonline{alenezi2016}; \citeonline{grano2019}; \citeonline{petrovic2021}; \citeonline{zhang2022} \\
    \hline
    Qualidade da seleção ou da aproximação do \textit{mutation score} & 7 & 19\% & \citeonline{zhang2010}; \citeonline{kurtz2016}; \citeonline{zhang2019}; \citeonline{kuroishi2020}; \citeonline{guizzo2022}; \citeonline{ashong2025}; \citeonline{ashong2026} \\
    \hline
    Custo e redução de mutantes (tempo de execução, número de mutantes, taxa de redução, esforço) & 6 & 17\% & \citeonline{just2012}; \citeonline{kurtz2016}; \citeonline{tufano2020}; \citeonline{guizzo2022}; \citeonline{arasteh2022}; \citeonline{arasteh2023} \\
    \hline
    Tamanho da suíte de testes e cobertura relativa entre critérios & 4 & 11\% & \citeonline{dimov2010}; \citeonline{kintis2010}; \citeonline{guerino2024}; \citeonline{ashong2026} \\
    \hline
    Validade frente a falhas reais (detecção de falhas reais, acoplamento) & 3 & 8\% & \citeonline{just2014}; \citeonline{petrovic2021}; \citeonline{khanfir2023} \\
    \hline
    Outras métricas isoladas (\textit{mutation brittleness}, \textit{test smells}, exposição a \textit{findings}) & 3 & 8\% & \citeonline{ojdanic2023}; \citeonline{grano2019}; \citeonline{petrovic2021} \\
\end{longtable}
\endgroup

O \textit{mutation score} é a razão entre os mutantes mortos e os mutantes não equivalentes \cite{chauhan2015, guerino2024}. Um mutante é morto quando algum teste produz resultado diferente do programa original; em suítes \textit{end-to-end} de aplicações web, quando ao menos um caso de teste falha, por uma ação ou por uma verificação (\textit{oracle}) \cite{yandrapally2021}. A métrica aparece com quatro variações. A primeira é o nome: \textit{kill score}, cobertura de mutação ou contagem de mutantes mortos \cite{inozemtseva2013, alenezi2016, dimov2010}. A segunda é o tratamento dos equivalentes: a maioria calcula o escore só sobre os não equivalentes, mas \citeonline{ahmed2016} e \citeonline{ashong2025} não os removem, pois a equivalência é indecidível \cite{yao2014}. A terceira é o recorte por operador ou só sobre mutantes dominadores, que elimina a redundância e evita uma percepção mais favorável da eficácia \cite{kurtz2016, kuroishi2020}. A quarta é o escore estimado sem executar todos os mutantes \cite{jalbert2012, zhang2019}. Há evidência de que o escore é um bom indicador: \citeonline{just2014} encontraram correlação significativa entre a detecção de mutantes e a de falhas reais, mais forte que a da cobertura de comandos.

A cobertura de código aparece em 9 estudos, sob a lógica de que falhas não são detectadas em código não coberto \cite{zhang2022}, e em sete deles acompanha o \textit{mutation score} como linha de base. Os estudos também mostram seus limites. A cobertura só indica que o código foi executado, e não que seu comportamento foi verificado \cite{petrovic2021}. Uma cobertura de linhas maior não implica suítes melhores: em \citeonline{alenezi2016}, a cobertura de mutação foi de 77\%, 61\% e 43\% em três projetos. \citeonline{ahmed2016} encontraram apenas correlação negativa fraca entre cobertura e correções de \textit{bugs} futuras, e \citeonline{assylbekov2013} reportam correlação de Pearson de 0,66 com o escore em um programa Java.

As demais métricas respondem a perguntas complementares. As contagens de mutantes explicam a sobrevivência: \citeonline{yao2014} mostram que a proporção de equivalentes e de teimosos varia por operador (o ABS gera muitos equivalentes e o LCR, muitos teimosos), e \citeonline{lin2018} usam as taxas de alcance, infecção e propagação. A qualidade da seleção de mutantes (7 estudos) usa o erro relativo do escore \cite{ashong2025} e o escore de mutantes dominadores \cite{kurtz2016}. O custo (6 estudos) é medido pelo tempo de execução, pelo número de mutantes e pela taxa de redução \cite{just2012, arasteh2022}. A validade frente a falhas reais (3 estudos) é medida pelo acoplamento: 70\% dos \textit{bugs} tinham ao menos um mutante acoplado em \citeonline{petrovic2021} e 73\% das falhas em \citeonline{just2014}.

Para aplicações web há uma única referência direta. \citeonline{yandrapally2021} propõem o MAEWU, que aplica 16 operadores ao DOM do navegador, e o avaliam em seis aplicações web de código aberto (\textit{AddressBook}, \textit{Claroline}, \textit{Collabtive}, \textit{MantisBT}, MRBS e PPMA), com suítes JUnit e Selenium. A métrica principal é o \textit{mutation score} sobre os não equivalentes, acrescido da gravidade do \textit{bug} e da teimosia (\textit{stubbornness}), uma escala de 0 a 4 do esforço para matar o mutante. Em 300 mutantes, 93\% eram não equivalentes, o escore médio foi de 0,20, e os mutantes de manipuladores de evento (EHAM, EHDM e EHMM) tiveram escore zero. Os outros estudos próximos tratam de composições WS-BPEL \cite{esterobotaro2010}, de ferramentas para Python \cite{guerino2024} e de código industrial \cite{petrovic2021}.

Em resumo, a literatura mede a eficácia de suítes sobretudo pelo \textit{mutation score} sobre mutantes não equivalentes, com a cobertura como linha de base. Isso sustenta adotá-lo como métrica principal da solução, com a cobertura como referência e a análise dos mutantes sobreviventes como apoio à interpretação.


\subsection[Quais operadores são mais utilizadas na literatura?]{Quais operadores de mutação e tipos de alterações artificiais no código-fonte são mais aplicados no contexto específico de aplicações web?}
\label{subsec:operadores-utilizadas}

Os operadores mais aplicados são os tradicionais, que alteram o código-fonte no nível de instruções e expressões. A maioria dos estudos (28 de 37) descreve os operadores ou o conjunto de operadores utilizado, e as famílias mais frequentes são a substituição de operadores relacionais, como o ROR (16 estudos, 44\%), a substituição de operadores aritméticos, como o AOR (15 estudos, 42\%), a substituição de operadores lógicos e condicionais, como o LCR e o COR (14 estudos, 39\%) e a remoção de instruções, blocos e chamadas de método, como o SDL e o SBR (14 estudos, 39\%). Em seguida aparecem os operadores unitários e de valor absoluto, como UOI e ABS, e a alteração de constantes, ambos com 11 estudos (31\%). Essas famílias coincidem com o conjunto seletivo clássico de cinco classes (ABS, AOR, LCR, ROR e UOI) avaliado por \citeonline{yao2014} e com os cinco operadores (AOR, LCR, ROR, UOI e SBR) do sistema de mutação do Google \cite{petrovic2021}.

Para aplicações web, apenas um estudo propõe operadores específicos: o MAEWU, com 16 operadores que alteram o DOM da página, e não o código-fonte \cite{yandrapally2021}. A Tabela~\ref{tab:familias-operadores} resume as famílias de alteração encontradas.

\begingroup
\footnotesize
\setlength{\tabcolsep}{4pt} % Ajusta o espaçamento interno lateral das células
\begin{longtable}{|p{0.20\textwidth}|p{0.18\textwidth}|>{\centering\arraybackslash}p{0.08\textwidth}|>{\centering\arraybackslash}p{0.06\textwidth}|p{0.40\textwidth}|}
    \caption{Famílias de operadores e tipos de alterações}
    \label{tab:familias-operadores} \\
    
    % --- Cabeçalho da PRIMEIRA página ---
    \hline
    \textbf{Família de alteração} & \textbf{Exemplos de operadores} & \textbf{Estudos} & \textbf{\%} & \textbf{Referências} \\
    \hline
    \endfirsthead

    % --- Cabeçalho das PÁGINAS SEGUINTES ---
    \multicolumn{5}{c}{\footnotesize\textit{Continuação da Tabela~\ref{tab:familias-operadores}}} \\
    \hline
    \textbf{Família de alteração} & \textbf{Exemplos de operadores} & \textbf{Estudos} & \textbf{\%} & \textbf{Referências} \\
    \hline
    \endhead

    % --- Rodapé intermediário ---
    \hline
    \endfoot

    % --- Rodapé FINAL ---
    \hline
    \multicolumn{5}{c}{\footnotesize Fonte: Autor (2026).} \\
    \endlastfoot

    % --- Conteúdo da tabela ---
    Substituição de operadores relacionais e de condições & ROR, Conditionals Boundary, Negate Conditionals & 16 & 44\% & \citeonline{just2012}; \citeonline{jalbert2012}; \citeonline{ojdanic2023}; \citeonline{belli2011}; \citeonline{tao2025}; \citeonline{petrovic2021}; \citeonline{zhang2022}; \citeonline{arasteh2023}; \citeonline{ashong2025a}; \citeonline{he2024}; \citeonline{kurtz2016}; \citeonline{just2014}; \citeonline{yao2014}; \citeonline{assylbekov2013}; \citeonline{ashong2025b}; \citeonline{inozemtseva2013} \\
    \hline
    Substituição de operadores aritméticos & AOR, AOR-B, AOR-S, AOR-U, Math, Increments & 15 & 42\% & \citeonline{jalbert2012}; \citeonline{esterobotaro2010}; \citeonline{belli2011}; \citeonline{guerino2024}; \citeonline{tao2025}; \citeonline{petrovic2021}; \citeonline{zhang2022}; \citeonline{arasteh2023}; \citeonline{ashong2025a}; \citeonline{he2024}; \citeonline{kurtz2016}; \citeonline{just2014}; \citeonline{yao2014}; \citeonline{assylbekov2013}; \citeonline{ashong2025b} \\
    \hline
    Substituição de operadores lógicos e condicionais & LCR, COR, LOR, COD, COI & 14 & 39\% & \citeonline{just2012}; \citeonline{esterobotaro2010}; \citeonline{belli2011}; \citeonline{guerino2024}; \citeonline{tao2025}; \citeonline{petrovic2021}; \citeonline{dimov2010}; \citeonline{arasteh2023}; \citeonline{ashong2025a}; \citeonline{he2024}; \citeonline{kurtz2016}; \citeonline{yao2014}; \citeonline{ashong2025b}; \citeonline{khanfir2023} \\
    \hline
    Remoção de instruções, blocos e chamadas de método & SDL, SBR, Remove Call, Void Method Call, RM & 14 & 39\% & \citeonline{jalbert2012}; \citeonline{chauhan2015}; \citeonline{ojdanic2023}; \citeonline{tao2025}; \citeonline{petrovic2021}; \citeonline{zhang2022}; \citeonline{ashong2025a}; \citeonline{he2024}; \citeonline{kurtz2016}; \citeonline{just2014}; \citeonline{assylbekov2013}; \citeonline{ashong2025b}; \citeonline{khanfir2023}; \citeonline{inozemtseva2013} \\
    \hline
    Inserção, remoção e substituição de operadores unários e valor absoluto & UOI, ABS, AOI, AOD, Invert Negatives & 11 & 31\% & \citeonline{jalbert2012}; \citeonline{ojdanic2023}; \citeonline{esterobotaro2010}; \citeonline{guerino2024}; \citeonline{tao2025}; \citeonline{petrovic2021}; \citeonline{zhang2022}; \citeonline{arasteh2023}; \citeonline{he2024}; \citeonline{yao2014}; \citeonline{assylbekov2013} \\
    \hline
    Substituição ou modificação de constantes e literais & CRP, Inline Constant, SRC & 11 & 31\% & \citeonline{jalbert2012}; \citeonline{chauhan2015}; \citeonline{ojdanic2023}; \citeonline{esterobotaro2010}; \citeonline{tao2025}; \citeonline{zhang2022}; \citeonline{he2024}; \citeonline{just2014}; \citeonline{assylbekov2013}; \citeonline{khanfir2023}; \citeonline{inozemtseva2013} \\
    \hline
    Alteração de retornos e de referências nulas & RSR, Return Values, NRV, NIV, NOI, RNC & 9 & 25\% & \citeonline{chauhan2015}; \citeonline{ojdanic2023}; \citeonline{tao2025}; \citeonline{zhang2022}; \citeonline{ashong2025a}; \citeonline{he2024}; \citeonline{assylbekov2013}; \citeonline{ashong2025b}; \citeonline{khanfir2023} \\
    \hline
    Substituição de variáveis & SVR, ISV, CSR & 6 & 17\% & \citeonline{jalbert2012}; \citeonline{chauhan2015}; \citeonline{esterobotaro2010}; \citeonline{tao2025}; \citeonline{he2024}; \citeonline{khanfir2023} \\
    \hline
    Operadores de deslocamento, de atribuição de atalho e bit a bit & SOR, SAOR & 5 & 14\% & \citeonline{guerino2024}; \citeonline{arasteh2023}; \citeonline{ashong2025a}; \citeonline{ashong2025b}; \citeonline{khanfir2023} \\
    \hline
    Operadores específicos de classe ou de domínio & Operadores de classe do MuJava, WS-BPEL, Python, DOM web (MAEWU) & 5 & 14\% & \citeonline{chauhan2015}; \citeonline{nica2011}; \citeonline{esterobotaro2010}; \citeonline{guerino2024}; \citeonline{yandrapally2021} \\
    \hline
    Geração alternativa (ordem superior, neural, baseada em \textit{bug reports}) & Mutação de segunda ordem, DeepMutation, iBiR & 3 & 8\% & \citeonline{kintis2010}; \citeonline{tufano2020}; \citeonline{khanfir2023} \\
\end{longtable}
\endgroup

Os operadores tradicionais se repetem entre ferramentas e linguagens. Os de expressões trocam um operador por outro do mesmo tipo (ROR, AOR, LCR e COR), como \texttt{\&\&} por \texttt{||} ou \texttt{>=} por \texttt{<=} \cite{belli2011}. Os de valores alteram constantes, inserem operador unário (UOI) ou valor absoluto (ABS) e substituem variáveis \cite{chauhan2015, yao2014}. Os de instruções removem código, como instruções (SDL), blocos (SBR) e chamadas de método, ou substituem o retorno \cite{jalbert2012, petrovic2021}. Os conjuntos aplicados vêm das ferramentas: o PIT oferece mutadores como \textit{Conditionals Boundary}, \textit{Negate Conditionals}, \textit{Math} e \textit{Return Values} \cite{assylbekov2013, zhang2022}; o Major, usado por \citeonline{just2014}, reúne a substituição de constantes e de operadores, a modificação de condições e a remoção de instruções; o LittleDarwin tem 14 operadores \cite{ashong2025a}; o Proteum implementa 108 operadores para C \cite{zhang2010}; e as ferramentas para Python trazem categorias de atribuição, aritméticos, unários, relacionais, condicionais e de \textit{bits} \cite{guerino2024}.

Nem todo operador se comporta da mesma forma. O COR e o ROR geram 41\% dos mutantes em \citeonline{just2012}, que redefiniram o COR sem redundância e reduziram em até 34\% o tempo da análise. No Google, o SBR gerou a maioria dos mutantes acoplados a falhas reais \cite{petrovic2021}. \citeonline{yao2014} mostram que o ABS e metade do UOI geram muitos equivalentes e quase nenhum teimoso, enquanto o LCR gera muitos teimosos.

Alguns estudos adaptam as alterações ao domínio ou mudam a forma de gerar os mutantes: operadores de classe do MuJava \cite{chauhan2015, nica2011}; 26 operadores para composições WS-BPEL, dos quais oito não foram aplicáveis às três composições avaliadas \cite{esterobotaro2010}; operadores específicos de Python, como alterações em decoradores \cite{guerino2024}; mutação de segunda ordem e mutação fraca \cite{kintis2010}; tradução neural, no DeepMutation \cite{tufano2020}; e injeção guiada por relatos de \textit{bugs}, no iBiR \cite{khanfir2023}.

Para aplicações \textit{web}, só um estudo propõe operadores próprios. \citeonline{yandrapally2021} observam que as ferramentas de mutação para a \textit{web} são ligadas a linguagens ou \textit{frameworks} específicos, como o \textit{WebMuJava}, que não é público, e o \textit{AjaxMutator}. O MAEWU altera o DOM dinâmico no navegador, sem acesso ao código-fonte, com 16 operadores derivados da análise de 250 relatos de \textit{bugs} e organizados em atributos, manipuladores de evento, árvore DOM, conteúdo e estilo. Dos 300 mutantes avaliados, 22 eram equivalentes (7 do SCM, cuja cor foi sobrescrita por uma regra CSS, e 5 do AAM), e os de manipuladores de evento foram os mais teimosos, com escore zero, pois matá-los exige executar a ação e verificar o resultado.

Em resumo, a literatura aplica sobretudo operadores tradicionais, como ROR, AOR, LCR e SDL, em programas Java, C e Python, e nenhum estudo da revisão os avaliou em suítes de interface \textit{web}. Para esse contexto, a única evidência usa operadores de DOM, que exigem uma suíte de interface. Esses resultados sustentam usar, na solução proposta, os operadores tradicionais, que são os mais frequentes, e deixar os operadores de DOM fora do escopo.


\subsection[Quais ferramentas são mais utilizadas na literatura?]{Quais ferramentas automatizadas de teste de mutação para tecnologias web são reportadas nos estudos primários e qual o seu nível de maturidade?}
\label{subsec:ferramentas-utilizadas}

Os estudos mencionam 28 ferramentas e sistemas de teste de mutação, usados nos experimentos ou citados como alternativas (Tabelas~\ref{tab:ferramentas-web} e~\ref{tab:ferramentas-proposito-geral}). Cinco se destinam a tecnologias \textit{web} (MAEWU, \textit{WebMuJava}, \textit{AjaxMutator}, GAmera e MutateMe), além de uma extensão do MuJava para componentes EJB, e somente o MAEWU foi concebido para suítes de testes de interface \textit{web} e está disponível publicamente \cite{yandrapally2021}. As demais são de propósito geral, voltadas sobretudo a Java, e são discutidas por representarem o estado da prática. Como os estudos não adotam um modelo de maturidade, as ferramentas foram classificadas por quatro evidências que os artigos reportam (disponibilidade pública, escopo de linguagens, adoção por outros estudos ou em produção e limitações relatadas) em cinco níveis: consolidada, estabelecida com ressalvas, ferramenta de pesquisa, indisponível ou inaplicável e não informada. Trata-se de uma síntese da autora, e não de um rótulo atribuído pelos estudos.

O MAEWU altera o DOM no navegador e só exige a URL da aplicação e a suíte de testes, sem acesso ao código nem uso de \textit{proxy}. Foi avaliado em seis aplicações com 50 mutantes cada, e seu repositório público traz instruções e dados rotulados \cite{yandrapally2021}. Com uma única avaliação, é classificado como ferramenta de pesquisa. Os autores não puderam usar as duas ferramentas anteriores: o \textit{WebMuJava}, para aplicações JSP, não é público, e o \textit{AjaxMutator}, para JavaScript do cliente, não analisou arquivos com mais de 50 mil linhas e não tinha mecanismo claro para aplicar os mutantes. Os demais estudos tratam de tecnologias específicas: o GAmera, para composições WS-BPEL \cite{esterobotaro2010}, uma extensão do MuJava para EJB \cite{nica2011} e o MutateMe, para PHP, sem descrição de maturidade \cite{belli2011}. Das seis ferramentas voltadas à \textit{web}, só o MAEWU está disponível e é aplicável a aplicações modernas, de modo que a maturidade das ferramentas para a \textit{web} é baixa.

\begin{table}[H]
    \centering
    \footnotesize
    \caption{Ferramentas de teste de mutação para tecnologias}
    \label{tab:ferramentas-web}
    \begin{tabular}{|p{2.3cm}|p{4.2cm}|p{4.8cm}|p{2.7cm}|}
        \hline
        \textbf{Ferramenta} & \textbf{Tecnologia-alvo e estudo} & \textbf{Evidência reportada} & \textbf{Maturidade} \\
        \hline
        MAEWU & Suítes de testes \textit{end-to-end} de aplicações web, por mutação do DOM dinâmico \cite{yandrapally2021} & Código público; 16 operadores; avaliada em seis aplicações de código aberto & Ferramenta de pesquisa \\
        \hline
        WebMuJava & Aplicações JSP e Java Server \cite{yandrapally2021} & Não está disponível publicamente & Indisponível \\
        \hline
        AjaxMutator & JavaScript do lado do cliente: Ajax e API do DOM \cite{yandrapally2021} & O Rhino não analisou arquivos com mais de 50 mil linhas; sem mecanismo claro para aplicar os mutantes & Inaplicável a aplicações modernas \\
        \hline
        GAmera & Composições WS-BPEL \cite{esterobotaro2010} & Gera e executa mutantes; avaliada em três composições & Ferramenta de pesquisa \\
        \hline
        MuJava estendido & Componentes EJB no JBoss \cite{nica2011} & Extensão para o JDK 1.5 e simulação de implantação; três aplicações EJB e um \textit{framework} da Siemens & Ferramenta de pesquisa \\
        \hline
        MutateMe & Código PHP \cite{belli2011} & Artigo não descreve funcionalidades nem maturidade & Não informada \\
        \hline
    \end{tabular}
    \legend{Fonte: Autor (2026).}
\end{table}

O PIT é a ferramenta de propósito geral mais usada (sete estudos), descrita como de ponta e capaz de alterar o \textit{bytecode} Java \cite{assylbekov2013, ahmed2016, zhang2022}. \citeonline{ashong2026} ressalvam que as ferramentas que alteram o código-fonte geram menos mutantes, mais fáceis de entender. O Major, integrado ao compilador Java, é usado em quatro estudos \cite{just2014, zhang2019, zhang2022}. O MuJava, usado em sete estudos, é bem estabelecido, mas descrito como desatualizado, e seu \textit{plugin} MuClipse não é mantido desde 2011 \cite{ramler2012, ashong2025a}. \citeonline{ramler2012} fazem a avaliação mais direta de maturidade: aplicam cinco ferramentas Java (Jumble, MuClipse, Simple Jester, Javalanche e PIT) a quatro suítes e concluem que atingiram um nível aceitável, com limitações técnicas, pois 10\% das execuções não produziram escore válido e a configuração exigiu \textit{scripts} de \textit{build}. O LittleDarwin e o MediumDarwin são ferramentas de pesquisa de um mesmo grupo.

\subsubsection{Ferramentas de propósito geral para Java}

O PIT é a ferramenta mais frequente: sete estudos o utilizam \cite{assylbekov2013, inozemtseva2013, ahmed2016, zhang2019, guizzo2022, zhang2022, ojdanic2023}. \citeonline{assylbekov2013} o descrevem como ferramenta de ponta para cobertura e teste de mutação em Java, que altera o \textit{bytecode} e executa os testes de unidade sobre as versões modificadas. \citeonline{ahmed2016} observam que o PIT já foi usado em vários estudos anteriores e o estendem para obter a matriz completa de falhas dos testes sobre os mutantes. \citeonline{zhang2022} o descrevem como amplamente usado e o empregam para gerar 163 conjuntos de dados originais e 650 adicionais. \citeonline{ashong2026} ressalvam que ele é mais rápido, mas que ferramentas que alteram o código-fonte geram menos mutantes, mais fáceis de entender para os profissionais. O PIT é classificado como \textbf{ferramenta consolidada}.

O Major é um \textit{framework} de mutação integrado ao compilador Java, que fornece o número de mutantes gerados e o tempo de execução da análise \cite{just2012}. \citeonline{just2014} o usam com operadores de substituição de constantes e de operadores, modificação de condições de desvio e remoção de instruções, e o listam, com o Javalanche e o PIT, entre as ferramentas Java modernas. \citeonline{zhang2022} o usam com seus oito operadores para gerar conjuntos de dados, e \citeonline{zhang2019} o empregam como segunda ferramenta, ao lado do PIT. \citeonline{alenezi2016} citam o Major e o MuJava, mas o uso experimental não está confirmado. O Major também é classificado como \textbf{ferramenta consolidada}.

O MuJava é usado diretamente em sete estudos \cite{kintis2010, nica2011, chauhan2015, shomali2020, arasteh2022, arasteh2023, tao2025}. Ele gera mutantes a partir do código-fonte, com operadores tradicionais adaptados à orientação a objetos e operadores de classe, executa os mutantes contra a suíte e reporta o \textit{mutation score} \cite{chauhan2015}. \citeonline{kintis2010} o descrevem como um sistema bem estabelecido e referenciado, e \citeonline{chauhan2015} e \citeonline{nica2011} o estendem com novos operadores e com suporte a EJB. Por outro lado, \citeonline{ashong2025} o chamam de desatualizado. \citeonline{ramler2012} acrescentam que o MuClipse, \textit{plugin} do Eclipse baseado no MuJava, não estava mais em desenvolvimento ativo desde 2011 e herdava a falta de suporte a \textit{generics} e anotações do MuJava. Por isso, o MuJava é classificado como \textbf{ferramenta estabelecida com ressalvas}.

O LittleDarwin e o MediumDarwin são \textbf{ferramentas de pesquisa}. O LittleDarwin gera mutantes alterando o código-fonte Java, ao contrário do PIT, o que reduz o desempenho, mas produz menos mutantes e mais fáceis de entender. Sua versão 0.10.9 suporta 14 operadores: nove padrão, que cobrem o conjunto clássico de Offutt, e cinco experimentais, que simulam falhas de tipo nulo e a remoção de métodos \cite{ashong2025}. O MediumDarwin v1.0.0.0 estende o LittleDarwin com análise de subsunção e com a lista completa de testes que matam cada mutante \cite{ashong2026}. O uso das duas ferramentas se concentra no mesmo grupo de pesquisa.

\citeonline{ramler2012} fazem a avaliação mais direta de maturidade na revisão. Eles aplicam cinco ferramentas Java (Jumble 1.1.0, MuClipse 1.3, Simple Jester 1.2, Javalanche 0.4 e PIT) a quatro suítes de testes de diferentes níveis de qualidade e concluem que as ferramentas atingiram um nível aceitável de maturidade, embora o uso prático ainda revele limitações técnicas. \textit{Plugins} para o Eclipse existiam para duas das cinco ferramentas, e as demais exigiram \textit{scripts} de linha de comando ou de \textit{build}, com maior esforço de configuração. O tempo médio de execução foi de 17,4 segundos, e 10\% das execuções (duas de vinte) não produziram um escore válido, no MuClipse e no Javalanche. O Jumble altera o \textit{bytecode} no nível de classe, com um conjunto limitado de operadores; o Simple Jester altera e recompila o código-fonte, o que aumenta o tempo de execução; e o Javalanche automatiza o processo por meio de um \textit{script} de \textit{build} cuja adaptação exigiu esforço considerável. \citeonline{arasteh2022} citam o Muclipse, o Jester e o Jumble como ferramentas compatíveis, sem usá-las, e \citeonline{jalbert2012} usam o Javalanche. Essas ferramentas são classificadas como \textbf{estabelecidas com ressalvas}.

\subsubsection{Ferramentas para outras linguagens, sistemas industriais e abordagens emergentes}

Para programas em C, a ferramenta usada é o Proteum, em três estudos \cite{zhang2010, kurtz2016, kuroishi2020}, na versão Proteum/IM. \citeonline{zhang2010} usam uma versão com 108 operadores de mutação, tradicionais e de interface. \citeonline{kurtz2016} a escolhem para os programas da suíte Siemens porque ela oferece um conjunto amplo de operadores em comparação com outras ferramentas, e citam o Mothra como origem dos operadores do conjunto E-seletivo. O Proteum é classificado como \textbf{ferramenta consolidada} para programas C. \citeonline{lin2018} apresentam o JCMuta, protótipo implementado em Java para programas ANSI-C, que integra 78 classes de operadores e gera mutantes fortes, fracos e de cobertura, e o descrevem como a primeira técnica capaz de medir a dificuldade de matar um mutante e de orientar a geração de testes para os mais difíceis. O Godzilla aparece apenas como referência, como ferramenta de geração de testes de mutação para programas Fortran \cite{kintis2010}.

Para Python, \citeonline{guerino2024} comparam, de forma estática e dinâmica, quatro ferramentas: Cosmic Ray, MutPy, MutMut e Mutatest. Os autores relatam a ausência de marcação de mutantes equivalentes e de redução de suítes, e apontam a execução de testes irrelevantes como oportunidade de otimização. \citeonline{tao2025} usam o MutPy 0.6.1 junto com o MuJava 2.0.

No ambiente industrial, \citeonline{petrovic2021} descrevem o sistema de teste de mutação do Google, implantado há mais de seis anos, integrado à revisão de código e com suporte a dez linguagens, entre elas JavaScript, TypeScript e Dart. O sistema já gerou mais de 15 milhões de mutantes e analisou centenas de milhares de \textit{changelists}, e reduz o custo mutando apenas o código alterado. Por ser interno, não pode ser reutilizado, mas é a evidência mais sólida de maturidade industrial entre os estudos.

As abordagens emergentes são \textbf{ferramentas de pesquisa}, quase todas restritas a Java. O DeepMutation, de \citeonline{tufano2020}, é a primeira ferramenta de mutação baseada em tradução neural, está disponível como código aberto e foi avaliada com projetos do Defects4J (\textit{Chart}, \textit{Lang}, \textit{Math} e \textit{Time}). Ela compilou com sucesso entre 35\% e 67\% dos projetos mutados e levou de 5 a 12 minutos para extrair, abstrair, mutar e traduzir os mutantes, sem contar a compilação. O iBiR, de \citeonline{khanfir2023}, injeta falhas guiadas por relatos de \textit{bugs}, invertendo os padrões de um sistema de reparo automático, e foi avaliado também com o Defects4J; o artigo cita o repositório público da ferramenta. O \textit{Predictive Mutation Testing}, de \citeonline{zhang2019}, prevê o resultado dos mutantes sem executá-los, é uma implementação experimental com página pública e deixa o serviço em nuvem como trabalho futuro. O Sentinel, de \citeonline{guizzo2022}, é uma hiper-heurística para gerar redutores de mutantes que utiliza o PIT 1.2.0. \citeonline{dimov2010} usam a \textit{Mutation Testing Platform} (MTP), ferramenta própria dos autores, sem informação sobre sua disponibilidade. Por fim, os geradores de testes EvoSuite, Randoop e JCrasher, usados por \citeonline{just2014}, e o EvoSuite, planejado por \citeonline{grano2019}, não são ferramentas de mutação, e dois estudos não reportam ferramentas de mutação: \citeonline{he2024}, sobre injeção de falhas, e \citeonline{yao2014}, que se baseiam em inspeção manual.

\begingroup
\footnotesize
\setlength{\tabcolsep}{4pt} % Ajusta o espaçamento interno lateral das células
\begin{longtable}{|p{2.7cm}|p{4.3cm}|p{4.3cm}|p{2.7cm}|}
    \caption{Ferramentas de teste de mutação de propósito geral}
    \label{tab:ferramentas-proposito-geral} \\
    
    % --- Cabeçalho da PRIMEIRA página ---
    \hline
    \textbf{Ferramenta} & \textbf{Alvo e estudos} & \textbf{Evidência reportada} & \textbf{Maturidade} \\
    \hline
    \endfirsthead

    % --- Cabeçalho das PÁGINAS SEGUINTES ---
    \multicolumn{4}{c}{\footnotesize\textit{Continuação da Tabela~\ref{tab:ferramentas-proposito-geral}}} \\
    \hline
    \textbf{Ferramenta} & \textbf{Alvo e estudos} & \textbf{Evidência reportada} & \textbf{Maturidade} \\
    \hline
    \endhead

    % --- Rodapé intermediário ---
    \hline
    \endfoot

    % --- Rodapé FINAL ---
    \hline
    \multicolumn{4}{c}{\footnotesize Fonte: Autor (2026).} \\
    \endlastfoot

    % --- Conteúdo da tabela ---
    PIT & Java, \textit{bytecode}; sete estudos \cite{assylbekov2013, inozemtseva2013, ahmed2016, zhang2019, guizzo2022, zhang2022, ojdanic2023} & Descrito como de ponta e amplamente usado; mais rápido que ferramentas de código-fonte & Consolidada \\
    \hline
    Major & Java, integrado ao compilador \cite{just2012, just2014, zhang2019, zhang2022, alenezi2016} & Listado entre as ferramentas Java modernas; uso confirmado em quatro estudos & Consolidada \\
    \hline
    MuJava & Java, código-fonte (sete estudos, de \citeonline{kintis2010} a \citeonline{tao2025}) & Bem estabelecido, mas descrito como desatualizado & Estabelecida com ressalvas \\
    \hline
    LittleDarwin & Java, código-fonte \cite{ashong2025} & Versão 0.10.9; 14 operadores & Ferramenta de pesquisa \\
    \hline
    MediumDarwin & Java, código-fonte \cite{ashong2026} & Versão 1.0.0.0; análise de subsunção & Ferramenta de pesquisa \\
    \hline
    Javalanche & Java, \textit{bytecode} \cite{jalbert2012, ramler2012} & \textit{Script} de \textit{build} de adaptação custosa; produziu escore inválido em um caso & Estabelecida com ressalvas \\
    \hline
    Jumble, MuClipse e Simple Jester (Jester) & Java \cite{ramler2012, arasteh2022} & Operadores limitados (Jumble); sem desenvolvimento ativo desde 2011 (MuClipse); tempos maiores (Simple Jester) & Estabelecida com ressalvas \\
    \hline
    Proteum (Proteum/IM) & C, 108 operadores \cite{zhang2010, kurtz2016, kuroishi2020} & Conjunto amplo de operadores & Consolidada \\
    \hline
    Mothra & Origem dos operadores E-seletivos \cite{kurtz2016} & Apenas citada & Não informada \\
    \hline
    JCMuta & ANSI-C, protótipo em Java, 78 classes de operadores \cite{lin2018} & Implementado pelos autores para a técnica proposta & Ferramenta de pesquisa \\
    \hline
    Godzilla & Fortran, geração de testes \cite{kintis2010} & Apenas citada & Não informada \\
    \hline
    Cosmic Ray, MutPy, MutMut e Mutatest & Python (\citeonline{guerino2024}; MutPy também em \citeonline{tao2025}) & Ausência de marcação de equivalentes e de redução de suítes & Não informada \\
    \hline
    Sistema do Google & Dez linguagens, incluindo JavaScript, TypeScript e Dart \cite{petrovic2021} & Mais de seis anos em produção; mais de 15 milhões de mutantes; interno & Consolidada (interna) \\
    \hline
    DeepMutation & Java, tradução neural \cite{tufano2020} & Código aberto; compila de 35\% a 67\% dos projetos mutados & Ferramenta de pesquisa \\
    \hline
    iBiR & Java, guiada por relatos de \textit{bugs} \cite{khanfir2023} & Avaliada com o Defects4J; repositório público & Ferramenta de pesquisa \\
    \hline
    \textit{Predictive Mutation Testing} & Java, previsão de resultados \cite{zhang2019} & Implementação experimental; serviço em nuvem como trabalho futuro & Ferramenta de pesquisa \\
    \hline
    Sentinel & Java, redutores de mutantes sobre o PIT 1.2.0 \cite{guizzo2022} & Hiper-heurística avaliada em sistemas de código aberto & Ferramenta de pesquisa \\
    \hline
    MTP & Plataforma própria para modelos de confiabilidade \cite{dimov2010} & Sem informação sobre disponibilidade & Não informada \\
\end{longtable}
\endgroup

Em síntese, a maturidade das ferramentas diminui à medida que se afasta de Java. As consolidadas são o PIT e o Major, para Java, o Proteum, para C, e o sistema do Google, e nenhuma delas foi concebida para a interface de aplicações \textit{web}. Para a \textit{web}, a única ferramenta disponível e aplicável é o MAEWU, uma ferramenta de pesquisa com uma única avaliação. Para Python, as quatro ferramentas comparadas por \citeonline{guerino2024} não têm maturidade informada nos estudos.

\subsection[Principais limitações e desafios práticos]{Quais são as principais limitações e desafios práticos identificados na aplicação de testes de mutação em aplicações web?}
\label{subsec:limitacoes-desafios}

Dos 37 estudos com dados extraídos, 33 relatam ao menos uma limitação ou desafio na aplicação do teste de mutação. Como somente quatro tratam de tecnologias web \cite{yandrapally2021, esterobotaro2010, nica2011, belli2011}, a maior parte das limitações vem de estudos sobre Java, C e Python. Elas são apresentadas aqui porque valem também para o contexto web, em especial o custo de execução e os mutantes equivalentes. Os desafios foram agrupados em seis temas (Tabela~\ref{tab:desafios-limitacoes}), e um mesmo estudo pode aparecer em mais de um deles. Os dois mais frequentes são o custo computacional, relatado em 22 estudos, e os mutantes equivalentes, em 20.

\begingroup
\footnotesize
\setlength{\tabcolsep}{4pt} % Ajusta o espaçamento interno lateral das células
\begin{longtable}{|p{2.2cm}|p{2.3cm}|p{3.2cm}|p{3.3cm}|}
    \caption{Principais limitações e desafios reportados nos estudos primários e estratégias de mitigação}
    \label{tab:desafios-limitacoes} \\
    
    % --- Cabeçalho da PRIMEIRA página ---
    \hline
    \textbf{Tema} & \textbf{Estudos (exemplos)} & \textbf{Limitação ou desafio} & \textbf{Estratégias reportadas} \\ 
    \hline
    \endfirsthead

    % --- Cabeçalho das PÁGINAS SEGUINTES ---
    \multicolumn{4}{c}{\footnotesize\textit{Continuação da Tabela~\ref{tab:desafios-limitacoes}}} \\
    \hline
    \textbf{Tema} & \textbf{Estudos (exemplos)} & \textbf{Limitação ou desafio} & \textbf{Estratégias reportadas} \\ 
    \hline
    \endhead

    % --- Rodapé intermediário ---
    \hline
    \endfoot

    % --- Rodapé FINAL ---
    \hline
    \multicolumn{4}{c}{\footnotesize Fonte: Autor (2026).} \\
    \endlastfoot

    % --- Conteúdo da tabela ---
    Especificidades de tecnologias web & \citeonline{yandrapally2021, esterobotaro2010, nica2011, belli2011} & Dados dinâmicos, fragilidade de localizadores, severidade de mutantes de estilo, serviços web indisponíveis, reimplantação a cada mutante & Mutação do DOM dinâmico, mockups de serviços, automação da reimplantação \\ 
    \hline
    Custo computacional & \citeonline{just2012, nica2011, petrovic2021, grano2019, guizzo2022} & Grande número de mutantes e de execuções de testes & Seleção de operadores, amostragem, mutação fraca, mutação de ordem superior, mutação apenas do código alterado, previsão do escore \\ 
    \hline
    Mutantes equivalentes & \citeonline{yao2014, he2024, ahmed2016, zhang2019} & Problema indecidível, com análise manual custosa e sujeita a erro & Revisão manual por mais de um avaliador, \textit{solver} de restrições, cálculo do escore sobre mutantes não equivalentes \\ 
    \hline
    Mutantes redundantes & \citeonline{ashong2026, kurtz2016, kuroishi2020, just2012} & Mutantes mortos pelos mesmos testes inflam o conjunto e distorcem o escore & Subsunção, mutantes dominadores, arcos primitivos \\ 
    \hline
    Representatividade dos mutantes e das métricas & \citeonline{just2014, tufano2020, khanfir2023, ashong2026} & Falhas reais sem mutante acoplado; mutantes pouco realistas; escore isolado & Novos operadores, mutação guiada por relatos de \textit{bugs}, seleção por histórico de defeitos \\ 
    \hline
    Ferramentas, adoção e validade dos estudos & \citeonline{ramler2012, petrovic2021, tufano2020, assylbekov2013} & Falhas das ferramentas, esforço de configuração, sobrecarga dos desenvolvedores, amostras pequenas & Regras de supressão, limite de mutantes por arquivo, \textit{corpora} variados e múltiplas repetições \\ 
\end{longtable}
\endgroup

Para a \textit{web}, \citeonline{yandrapally2021} apontam os dados dinâmicos, como valores de \textit{id} gerados em tempo de execução, cujos mutantes se parecem com equivalentes: dos 300 mutantes, 22 eram equivalentes, 7 do operador SCM (a cor era sobrescrita por uma regra CSS) e 5 do AAM (atributos não usados pelo JavaScript). Também citam a fragilidade dos localizadores de elementos \textit{web}, que encarece a manutenção das suítes \textit{end-to-end}, um modelo de severidade que trata o CSS como cosmético e subestima os mutantes de estilo, e a amostra limitada a seis aplicações, com rótulos atribuídos manualmente pelo primeiro autor. Para serviços e componentes \textit{web}, os desafios são de infraestrutura. Em WS-BPEL, os serviços reais não estavam disponíveis, os \textit{mockups} aumentam os mutantes potencialmente equivalentes e a implantação de cada mutante, com \textit{timeouts}, é demorada \cite{esterobotaro2010}. Em EJB, cada alteração exige nova implantação no servidor, e o experimento foi encerrado pelo volume de tempo e de recursos \cite{nica2011}. \citeonline{belli2011} apontam sobreviventes que exigem estados específicos do sistema e a aplicabilidade dos operadores ao PHP.

Os mesmos autores destacam que um dos maiores desafios de manter suítes de testes \textit{end-to-end} é a fragilidade dos localizadores de elementos web, que exige análise manual e manutenção custosas \cite{yandrapally2021}. Eles observam ainda que o modelo de severidade de \textit{bugs} usado na avaliação se baseia em aplicações mais antigas e trata problemas de CSS como cosméticos, o que subestima mutantes de estilo (SCM, SSM e SPM) em aplicações modernas com \textit{layouts} fluidos. Por fim, reconhecem amostra limitada, com seis aplicações e 300 mutantes, e rótulos atribuídos manualmente pelo primeiro autor, sem uma referência automatizada. As ferramentas anteriores também tiveram limitações práticas: o WebMuJava não está disponível e o AjaxMutator não analisou o JavaScript das aplicações modernas, conforme descrito na Seção~\ref{subsec:ferramentas-utilizadas}.

Os demais estudos sobre tecnologias web mostram desafios ligados à implantação e à infraestrutura. \citeonline{esterobotaro2010} relatam que oito operadores para composições WS-BPEL não se aplicaram a nenhuma das três composições avaliadas e que os serviços web reais não estavam disponíveis. Os \textit{mockups} usados no lugar deles não refletem a complexidade dos serviços reais e tendem a aumentar o número de mutantes potencialmente equivalentes. Os autores acrescentam que o tempo de implantação e remoção de cada mutante, além de \textit{timeouts} na troca de mensagens, pode ser significativo. \citeonline{nica2011} descreve o mesmo problema em componentes EJB: cada alteração na lógica de negócio exige uma nova implantação no servidor de aplicação, o que, para muitos mutantes, aumenta o tempo e o custo, e o servidor pode sofrer gargalos. O experimento com EJB foi temporariamente encerrado pelo volume de dados e de tempo e pela falta de recursos financeiros. Como mitigação, a autora propõe automatizar a reimplantação e a execução dos testes e injetar falhas apenas nos métodos de negócio que interferem na comunicação entre cliente e servidor. \citeonline{belli2011} apontam dificuldades com mutantes sobreviventes que exigem estados específicos do sistema ou extensões de cobertura consideradas impraticáveis, e com a aplicabilidade dos operadores ao código PHP.

\subsubsection{Custo computacional e de execução}
\label{subsubsec:custo-computacional}

O custo é a limitação mais citada, e aparece em 22 estudos. \citeonline{petrovic2021} consideram a abordagem tradicional, que avalia todos os mutantes, computacionalmente inviável na escala do Google; até o \textit{mutation score} é difícil de transformar em ação pelos desenvolvedores. A geração de mutantes também não é trivial e carrega um custo oculto, como indexar as unidades de compilação de C++, e a análise de acoplamento envolveu quase 400 mil mutantes e mais de 33 milhões de execuções de alvos de teste. Como não é possível avaliar todos os alvos de teste, os mutantes vivos reportados ficam imprecisos e a confiança dos desenvolvedores diminui. \citeonline{grano2019} descreve o custo como extremo, e \citeonline{arasteh2022, arasteh2023} o tratam como problema estrutural da técnica, em tempo e em dinheiro.

Outros estudos detalham o custo por ângulos diferentes. \citeonline{jalbert2012} o tomam como motivação central para prever o escore sem executar todos os mutantes. \citeonline{guizzo2022} observam que reduzir o número de mutantes não garante redução proporcional do tempo, e por isso medem o tempo de execução diretamente. \citeonline{inozemtseva2013} notam que paralelizar a execução acelera o processo, mas não elimina o problema da quantidade de mutantes. \citeonline{ojdanic2023} apontam o custo de aplicação em integração contínua e de cálculo da subsunção, e a dificuldade de comparar escores quando o conjunto de mutantes é recriado a cada versão. \citeonline{kintis2010} apontam que o número de mutantes candidatos é enorme, estimado na ordem do quadrado do número de linhas do programa. \citeonline{ashong2025} usam 15 projetos Java, número que, segundo os dados extraídos, foi restringido pelo alto custo da análise de mutação.

\subsubsection{Mutantes equivalentes e redundantes}
\label{subsubsec:mutantes-equivalentes}

A identificação de mutantes equivalentes é relatada em 20 estudos como limitação. O problema é indecidível, e \citeonline{yao2014} usaram análise manual, que consumiu seis meses-pessoa de esforço contínuo para classificar 1.230 mutantes de 18 programas. Segundo \citeonline{he2024}, a proporção de equivalentes entre os mutantes gerados costuma variar de 10\% a 40\%, o que dificulta julgar a adequação dos testes. \citeonline{zhang2019} reconhecem que não é possível garantir uma identificação automática geral, o que atrapalha o cálculo de um escore que os exclua corretamente. Na prática, muitos estudos assumem o risco: \citeonline{ahmed2016} não aplicaram método sistemático de remoção e admitem impacto no escore, e \citeonline{ashong2025} tratam todos os sobreviventes como matáveis ($E = 0$). \citeonline{guerino2024} observam que as ferramentas Python avaliadas não marcam mutantes equivalentes e, por isso, adotam revisão manual por um segundo avaliador. \citeonline{nica2011} propõe, ainda em investigação, presumir equivalentes quando um \textit{solver} de restrições não encontra um teste que os distinga.

Os mutantes redundantes são outro desafio, pois não são equivalentes, mas são mortos pelos mesmos testes e inflam o conjunto sem aumentar a precisão \cite{ashong2026}. \citeonline{kurtz2016} usam mutantes dominadores justamente para eliminar o ruído causado pela redundância, mas notam que o problema geral de encontrar conjuntos de dominadores é indecidível. \citeonline{kuroishi2020} apontam que o próprio cálculo do conjunto mínimo, usado para avaliar a redundância, exige processamento. \citeonline{tao2025} identificam ainda redundância nos dados de teste gerados, porque uma mesma entrada pode cobrir vários ramos de mutação.

\subsubsection{Representatividade dos mutantes e das métricas}
\label{subsubsec:representatividade}

Uma limitação de fundo é que os mutantes podem não representar as falhas reais. \citeonline{just2014} mostram que 27\% das falhas reais não têm mutante acoplado entre os gerados por operadores comuns e que 17\% não têm mutante acoplado de nenhum tipo, em geral por envolverem modificação ou simplificação de algoritmo (37 casos) ou remoção de código (7 casos). \citeonline{tufano2020} tratam a representatividade como desafio central, pois mutantes pouco realistas podem subestimar a eficácia de uma suíte, e as taxonomias de mutantes costumam exigir grande esforço para serem criadas. \citeonline{khanfir2023} acrescentam que o problema envolve tanto o tipo da alteração quanto o local em que ela é inserida. \citeonline{inozemtseva2013} alertam que a redução aleatória e a redução por operadores podem eliminar mutantes informativos. \citeonline{kurtz2016} concluem que uma seleção de operadores muito eficaz em um programa tende a ter desempenho ruim em outro.

Há também limitações nas métricas. \citeonline{ashong2026} afirmam que o \textit{mutation score} isolado não capta a diversidade e a capacidade de revelar falhas do conjunto de mutantes, e por isso avaliam as técnicas sob quatro métricas. \citeonline{grano2019} aponta que a cobertura de código também é um indicador não ótimo da qualidade da suíte.

\subsubsection{Limitações das ferramentas e da adoção prática}
\label{subsubsec:limitacoes-ferramentas}

\citeonline{ramler2012} observam que, mesmo em cenários simples, 10\% das execuções das cinco ferramentas Java não produziram um escore válido, e que as ferramentas sem \textit{plugin} exigiram \textit{scripts} de linha de comando ou de \textit{build}, com configuração mais complexa. \citeonline{tufano2020} relatam que apenas 35\% a 67\% dos projetos mutados pelo DeepMutation compilaram, que a ferramenta opera só no nível de método e que, na avaliação, métodos com mais de 50 \textit{tokens} foram descartados. \citeonline{guerino2024} apontam, nas ferramentas Python, a ausência de redução de suítes e a execução de testes irrelevantes. \citeonline{assylbekov2013} registram que a grande variedade de ferramentas dificulta uma medição padronizada e que ferramentas diferentes podem alterar levemente a correlação observada. \citeonline{nica2011} limita-se a mutação de código-fonte no nível de teste de unidade, com a geração de testes distinguidores avaliada em programas de até 100 linhas.

Sobre a adoção por equipes de desenvolvimento, \citeonline{petrovic2021} mostram que o Google precisou adaptar a técnica: mutar apenas o código alterado e coberto por testes, gerar um mutante por linha, aplicar regras de supressão para código que não produz mutantes produtivos, como \textit{logging}, e relatar no máximo sete mutantes por arquivo para evitar sobrecarga cognitiva. A abordagem de \citeonline{arasteh2023} depende da calibração de parâmetros do algoritmo de otimização para cada programa, e \citeonline{lin2018} reconhecem que a análise dinâmica de mutantes teimosos exige uma grande quantidade de testes, o que a torna impraticável antes de a suíte existir.

\subsubsection{Limitações de validade dos estudos}
\label{subsubsec:limitacoes-validade}

Vários estudos avaliam suas técnicas em amostras pequenas. \citeonline{assylbekov2013} usam programas Java de no máximo 260 linhas, \citeonline{lin2018} usam programas com menos de mil linhas e observam que, com o aumento da complexidade, o número de entradas e de subdomínios cresce de forma explosiva e encarece a estimativa. \citeonline{kurtz2016} reconhecem que os programas da suíte Siemens, embora muito usados, são pequenos e podem não ser representativos de programas reais. \citeonline{yandrapally2021} avaliam seis aplicações web, e \citeonline{ashong2025}, 15 projetos Java. \citeonline{assylbekov2013} também notam que a qualidade dos testes depende da proficiência do programador e que a seleção de casos de teste não foi aleatória. Como mitigação, os estudos recorrem a \textit{corpora} de diferentes tamanhos e domínios, repetições múltiplas das técnicas estocásticas (1.000 por configuração em \citeonline{ashong2025, ashong2026}) e disponibilização pública de ferramentas e dados, como \citeonline{yandrapally2021}.

\subsubsection{Síntese e implicações}
\label{subsubsec:sintese-implicacoes}

Em síntese, o custo de execução e os mutantes equivalentes são limitações quase universais, presentes em estudos de todas as linguagens e ferramentas, e serão enfrentados em qualquer aplicação web. Para a web, os estudos acrescentam obstáculos próprios: dados dinâmicos que geram mutantes quase equivalentes, fragilidade dos localizadores das suítes de interface, reimplantação a cada mutante em servidores de aplicação, dependência de serviços externos e ferramentas indisponíveis ou inaplicáveis. A evidência sobre esses obstáculos é pequena, pois se apoia em quatro estudos e em amostras reduzidas. Esses resultados sugerem que a solução proposta deva limitar o número de mutantes por aplicação, como faz o MAEWU com 50 mutantes, priorizar código alterado ou coberto, tratar equivalentes com revisão manual por mais de um avaliador e reportar o \textit{mutation score} junto de métricas complementares.

\subsection[Estratégias para reduzir os custos]{Quais estratégias são propostas na literatura para reduzir os custos e mitigar as limitações da aplicação de testes de mutação em aplicações web?}
\label{subsec:estrategias-reducao}

Dos 37 estudos com dados extraídos, 31 propõem ou avaliam ao menos uma estratégia para reduzir o custo ou mitigar limitações do teste de mutação. Quase todas foram desenvolvidas e avaliadas em programas Java, C ou Python, e só quatro estudos tratam de tecnologias web: o MAEWU \cite{yandrapally2021}, os compostos WS-BPEL \cite{esterobotaro2010}, os componentes EJB \cite{nica2011} e o código PHP \cite{belli2011}. Por isso, o tópico apresenta primeiro as estratégias feitas para a web e, depois, as demais, indicando o que cada uma poderia oferecer a aplicações web. \citeonline{grano2019} lembra que Offutt e Untch classificam as técnicas de aceleração do teste de mutação em três grupos: fazer menos (\textit{do fewer}), fazer de forma mais inteligente (\textit{do smarter}) e fazer mais rápido (\textit{do faster}). A maioria das estratégias encontradas se enquadra no primeiro grupo, ao reduzir o número de mutantes (Tabela~\ref{tab:estrategias-reducao}).

\begin{table}[H]
    \centering
    \footnotesize
    \caption{Estratégias para reduzir custos e mitigar limitações}
    \label{tab:estrategias-reducao}
    \begin{tabular}{|p{3.1cm}|p{3.8cm}|p{4.6cm}|p{2.5cm}|}
        \hline
        \textbf{Estratégia} & \textbf{Estudos (exemplos)} & \textbf{Resultado ou evidência reportada} & \textbf{Uso em web} \\
        \hline
        Seleção por importância, com mutantes espalhados (MAEWU) & \citeonline{yandrapally2021} & Limite de 50 mutantes por aplicação, escolhidos por características estáticas e dinâmicas dos elementos web & Avaliada em seis aplicações web \\
        \hline
        Infraestrutura e testes para serviços e componentes web & \citeonline{esterobotaro2010}; \citeonline{nica2011} & \textit{Mockups} de serviços, suíte mínima (28 para 7 testes), reimplantação automatizada, mutação só dos métodos de negócio & WS-BPEL e EJB \\
        \hline
        Seleção de operadores e remoção de operadores redundantes & \citeonline{zhang2010}; \citeonline{kurtz2016}; \citeonline{yao2014}; \citeonline{just2012} & Seleção de operadores não superou a aleatória; remoção de redundantes reduziu até 34\% do tempo & Não avaliada \\
        \hline
        Amostragem e seleção estática de mutantes & \citeonline{ashong2025, ashong2026}; \citeonline{inozemtseva2013} & Nenhuma técnica é a melhor em todos os casos; as estratificadas se destacam & Não avaliada \\
        \hline
        Mutação em caminhos propensos a falhas & \citeonline{arasteh2022, arasteh2023}; \citeonline{shomali2020} & Redução de 28,10\% e de 27,63\% dos mutantes & Não avaliada \\
        \hline
        Mutantes dominadores, subsunção e arcos primitivos & \citeonline{kurtz2016}; \citeonline{kuroishi2020}; \citeonline{ojdanic2023} & Redução da redundância e do ruído no escore & Não avaliada \\
        \hline
        Mutação de ordem superior e mutação fraca & \citeonline{kintis2010}; \citeonline{tao2025} & 96\% a 98\% da cobertura de mutação forte, com menos equivalentes & Não avaliada \\
        \hline
        Previsão do escore sem executar os mutantes & \citeonline{jalbert2012}; \citeonline{zhang2019}; \citeonline{zhang2022}; \citeonline{grano2019} & Modelos supervisionados e abordagem não supervisionada baseada em cobertura & Não avaliada \\
        \hline
        Mutação orientada a mudanças e a regras de contexto & \citeonline{petrovic2021}; \citeonline{ojdanic2023} & Aplicada em produção no Google, com suporte a JavaScript, TypeScript e Dart & Parcial (linguagens web) \\
        \hline
        Mutantes realistas e direcionados & \citeonline{tufano2020}; \citeonline{khanfir2023}; \citeonline{just2014}; \citeonline{guizzo2022} & iBiR gera falhas acopladas a falhas reais em cerca de 37\% dos casos, contra menos de 4\% da mutação tradicional & Não avaliada \\
        \hline
        Tratamento de mutantes equivalentes & \citeonline{nica2011}; \citeonline{guerino2024}; \citeonline{yao2014}; \citeonline{he2024} & Revisão manual, \textit{solver} de restrições, modelos de geração de testes por LLM & Não avaliada \\
        \hline
    \end{tabular}
    \legend{Fonte: Autor (2026).}
\end{table}

A estratégia mais relevante para este trabalho é a do MAEWU, que trata o custo de mutar páginas \textit{web} pesadas: uma página média tem cerca de 2~MB e mais de 600 elementos de 32 tipos, e sem uma noção de importância o número de mutantes se torna ingerenciável \cite{yandrapally2021}. A seleção é aleatória, enviesada e espalhada: a chance de escolher um mutante depende de uma pontuação que soma características estáticas do elemento (ser folha, ter texto, ser interativo e estar visível) e dinâmicas, do traço de execução (acessos pelos testes e pelo JavaScript), e que cai quando já há mutantes no mesmo elemento. O limite foi de 50 mutantes por aplicação. Para serviços e componentes \textit{web}, \citeonline{esterobotaro2010} usam \textit{mockups} e suítes mínimas e não redundantes (28 para 7 testes em uma composição), e \citeonline{nica2011} propõe mutar só os métodos de negócio que interferem na comunicação entre cliente e servidor e automatizar a reimplantação. Todas foram avaliadas em pequena escala e dependem da tecnologia.

Para reduzir o número de mutantes, a seleção de operadores é a estratégia mais tradicional, mas não supera a aleatória \cite{zhang2010} e tem desempenho ruim em ao menos alguns programas \cite{kurtz2016}. \citeonline{just2012} mostram que o COR e o ROR, que criam 41\% dos mutantes, podem ser redefinidos sem redundância, o que reduz o tempo em até 34\%. A amostragem e a seleção estática são mais gerais: \citeonline{ashong2025a} comparam dez técnicas em 15 projetos Java e concluem que nenhuma é a melhor em todos os casos, mas que as de estratificação superam as demais em até 129 dos 135 cenários, e que a estratificação pelo arquivo de origem deu a melhor aproximação do escore em até 68. Outra linha concentra a mutação em caminhos propensos a falhas, com reduções de 28,10\% e de 27,63\% dos mutantes \cite{arasteh2022, arasteh2023}. Para a redundância, usam-se mutantes dominadores e subsunção \cite{kurtz2016, ojdanic2023}, e a mutação de ordem superior e a mutação fraca atingem de 96\% a 98\% da cobertura de mutação forte com menos equivalentes \cite{kintis2010}.

Para serviços e componentes web, \citeonline{esterobotaro2010} recorrem a \textit{mockups} para contornar a indisponibilidade dos serviços reais e a suítes de testes adequadas, não redundantes e mínimas, para que a eficácia dos operadores não seja distorcida por testes improdutivos. Em um dos casos, a suíte da composição \textit{MetaSearch} foi reduzida de 28 para 7 casos de teste. \citeonline{nica2011} propõe, para EJB, definir um conjunto eficaz de operadores, escolher um conjunto de testes que revele a maioria das falhas, mutar apenas os métodos de negócio que interferem na comunicação entre cliente e servidor e automatizar a reimplantação e a execução dos testes. Também propõe gerar casos de teste distinguidores com um \textit{solver} de restrições, que matam mutantes vivos sem gerar testes redundantes. Todas essas estratégias foram avaliadas em escala pequena e dependem de decisões específicas de cada tecnologia.

\subsubsection{Seleção e redução do número de mutantes}
\label{subsubsec:selecao-reducao}

A seleção de operadores é a estratégia mais tradicional, mas os estudos questionam sua eficácia. \citeonline{zhang2010} comparam três técnicas de seleção por operadores com duas aleatórias e concluem que a seleção por operadores não é superior à aleatória, que pode ser até melhor. \citeonline{kurtz2016} avaliam, no Proteum, todos os pequenos conjuntos de operadores e mostram que todas as abordagens de seleção têm desempenho ruim em ao menos alguns programas, o que sugere uma seleção especializada, escolhida conforme o artefato testado. \citeonline{yao2014} recomendam priorizar operadores que geram muitos mutantes teimosos e poucos equivalentes, como o LCR, e evitar a classe ABS e metade do UOI, que geram muitos equivalentes. \citeonline{just2012} mostram que os operadores COR e ROR, que criam 41\% dos mutantes nos programas estudados, podem ser redefinidos sem redundância, o que reduz o tempo da análise em até 34\% e evita um escore superestimado em até 10\%.

A amostragem e a seleção estática de mutantes são mais gerais. \citeonline{ashong2025} comparam dez técnicas estáticas de seleção em 15 projetos Java e concluem que nenhuma é a melhor em todos os projetos e taxas de amostragem, mas que as técnicas baseadas em estratificação superam as demais em até 129 dos 135 cenários; a estratificação pelo arquivo de origem deu a melhor aproximação do escore em quase metade dos casos (até 68 dos 135). \citeonline{ashong2026} estendem essa avaliação com quatro métricas, para considerar também a diversidade dos mutantes selecionados. \citeonline{inozemtseva2013} selecionam apenas mutantes de arquivos com muitos defeitos históricos, usando a seleção aleatória, com 776 mutantes, como referência. \citeonline{guizzo2022} propõem o Sentinel, uma hiper-heurística baseada em evolução gramatical e otimização multiobjetivo que gera automaticamente estratégias de redução, com \textit{cache} e reutilização das estratégias entre versões.

Em síntese, as estratégias mais estudadas reduzem o número de mutantes, e nenhuma técnica é superior em todos os contextos. Para a \textit{web}, a única estratégia de seleção avaliada é a do MAEWU, e as demais dependem de WS-BPEL e EJB. As estratégias gerais de Java, C e Python não foram avaliadas em aplicações \textit{web}, mas a mutação orientada a mudanças do Google e a previsão do escore são candidatas naturais por não dependerem da linguagem. Esses resultados sustentam combinar a amostragem estratificada com a mutação restrita ao código coberto ou alterado.


\subsection{Quais fatores são associados à sobrevivência de mutantes nas suítes de testes avaliadas?}
\label{subsec:fatores-sobrevivencia}

Dos 37 estudos com dados extraídos, 18 reportam ao menos um fator associado à sobrevivência de mutantes, um trata os sobreviventes como equivalentes sem analisá-los \cite{khanfir2023} e 17 não informam fatores. Os fatores foram agrupados em oito categorias, apresentadas na Tabela~\ref{tab:fatores-sobrevivencia}, e um mesmo estudo pode contribuir para mais de uma delas. As evidências vêm predominantemente de programas Java e C avaliados com testes de unidade, enquanto apenas \citeonline{yandrapally2021} investigam a sobrevivência de mutantes em suítes de testes de interface de aplicações web.

\begin{table}[H]
    \centering
    \footnotesize
    \caption{Fatores associados à sobrevivência de mutantes}
    \label{tab:fatores-sobrevivencia}
    \begin{tabular}{|p{3.0cm}|p{3.8cm}|c|c|p{4.3cm}|}
        \hline
        \textbf{Fator} & \textbf{Exemplos} & \textbf{Estudos} & \textbf{\%} & \textbf{Referências} \\
        \hline
        Características da suíte de testes & Qualidade e forma de construção dos testes, oráculos, utilidade marginal decrescente & 10 & 28\% & \cite{esterobotaro2010, assylbekov2013, inozemtseva2013, alenezi2016, lin2018, grano2019, petrovic2021, yandrapally2021, zhang2022, ashong2026} \\
        \hline
        Equivalência & Mutante com o mesmo comportamento do original; trechos sem efeito observável & 9 & 25\% & \cite{zhang2019, belli2011, tao2025, yandrapally2021, he2024, just2014, yao2014, assylbekov2013, inozemtseva2013} \\
        \hline
        Tipo de operador e dificuldade de detecção & Mutantes teimosos ou resistentes; operadores que exigem testes específicos & 7 & 19\% & \cite{esterobotaro2010, belli2011, lin2018, yao2014, zhang2019, tao2025, yandrapally2021} \\
        \hline
        Características e importância do código mutado & Tamanho e complexidade; funcionalidades não implementadas; histórico de defeitos; uso pelos usuários & 5 & 14\% & \cite{belli2011, inozemtseva2013, yao2014, ahmed2016, grano2019} \\
        \hline
        Cobertura do código mutado & Local da mutação não executado; poucos testes cobrindo o mutante & 4 & 11\% & \cite{just2014, zhang2019, grano2019, zhang2022} \\
        \hline
        Falha de alcance, infecção ou propagação do estado & Estado não infectado; efeito não observável na saída & 3 & 8\% & \cite{yao2014, lin2018, zhang2019} \\
        \hline
        Redundância entre mutantes e estratégia de mutação & Mutantes que compartilham o mesmo destino; mutação fraca; mutantes disjuntos & 2 & 6\% & \cite{kintis2010, petrovic2021} \\
        \hline
        Falhas reais não representadas pelos mutantes & Operadores fracos, ausentes ou inexistentes; falhas de algoritmo e remoção de código & 2 & 6\% & \cite{just2014, petrovic2021} \\
        \hline
    \end{tabular}
    \legend{Fonte: Autor (2026).}
\end{table}

\subsubsection{Características da suíte de testes}
\label{subsubsec:caracteristicas-suite}

Esta é a categoria mais frequente, identificada em 10 estudos (28\%). A sobrevivência está relacionada principalmente à qualidade dos casos de teste, às verificações realizadas e à forma como a suíte foi construída. \citeonline{alenezi2016} demonstram que uma quantidade maior de testes não implica necessariamente melhor cobertura de mutação, enquanto \citeonline{grano2019} associam maior eficácia à cobertura de comandos e à ausência de \textit{test smells}.

Além disso, \citeonline{assylbekov2013} apontam que testes baseados em entradas aleatórias podem deixar mutantes sobreviventes, enquanto \citeonline{zhang2022} discutem a utilidade marginal decrescente dos testes, na qual a adição de novos casos tende a produzir ganhos progressivamente menores na detecção de mutantes.

\subsubsection{Equivalência}
\label{subsubsec:equivalencia}

A equivalência aparece em nove estudos (25\%). Mutantes equivalentes apresentam o mesmo comportamento do programa original e, portanto, não podem ser mortos por nenhum caso de teste \cite{assylbekov2013, tao2025}. Sua presença dificulta a interpretação da sobrevivência, pois nem todo mutante sobrevivente representa uma deficiência da suíte.

\citeonline{yao2014} identificam que a equivalência pode ocorrer quando a alteração não modifica o estado do programa ou quando seus efeitos não alcançam uma saída observável. Os autores também demonstram que a frequência de mutantes equivalentes varia conforme o operador utilizado, reforçando a necessidade de considerar suas características durante a análise.

\subsubsection{Tipo de operador e dificuldade de detecção}
\label{subsubsec:tipo-operador}

O tipo de operador e a dificuldade de detecção foram identificados em sete estudos (19\%). Determinados operadores produzem mutantes não equivalentes que exigem entradas ou verificações específicas para serem mortos, frequentemente denominados mutantes teimosos (\textit{stubborn mutants}) \cite{yao2014, zhang2019}.

\citeonline{tao2025} associam essa dificuldade à pequena quantidade de entradas capazes de detectar determinadas alterações. Já \citeonline{lin2018} demonstram que a sobrevivência pode depender tanto das características dos mutantes quanto das estratégias utilizadas para selecionar os casos de teste. Dessa forma, diferentes operadores podem apresentar níveis distintos de dificuldade de detecção.

\subsubsection{Demais fatores e síntese dos resultados}
\label{subsubsec:sintese-fatores}

Os demais fatores incluem características do código mutado (5 estudos), cobertura do código (4), falhas de alcance, infecção ou propagação do estado (3), redundância entre mutantes (2) e limitações na representatividade de defeitos reais (2). Esses estudos indicam que a sobrevivência também pode depender da complexidade do programa, da ausência de execução do trecho modificado e da impossibilidade de observar seus efeitos \cite{grano2019, just2014, lin2018, zhang2019}.

Em síntese, a literatura associa a sobrevivência principalmente às características da suíte de testes, à equivalência dos mutantes e à dificuldade de detecção dos operadores. Esses fatores mostram que um mutante sobrevivente não indica necessariamente uma deficiência dos testes, sendo importante distinguir alterações equivalentes daquelas que poderiam ser detectadas por testes mais adequados. A escassez de evidências específicas sobre aplicações web também reforça a necessidade de investigar esses fatores nesse contexto.

\subsection[Distribuição por ano sobre testes de mutação]{Qual é a distribuição, por ano, das publicações sobre testes de mutação aplicados à avaliação de suítes de testes automatizados em aplicações web?}
\label{subsec:distribucao-ano}

Os 37 estudos com dados extraídos foram publicados entre 2010 e 2026, com ao menos um estudo em 16 dos 17 anos do intervalo; só 2017 não tem nenhum. O ano vem do campo ``Ano de Publicação'' da planilha de extração. A Figura~\ref{fig:publicacoes-ano-web} e a Tabela~\ref{tab:distribuicao-publicacoes} mostram a distribuição. A maior parte dos estudos não trata de aplicações web: apenas \citeonline{yandrapally2021} avaliam suítes de testes de aplicações web como objeto principal, e outros dois estudos têm relação indireta com o domínio web. Por isso, a distribuição é discutida a partir de todos os estudos, e a Seção~\ref{subsubsec:publicacoes-web} trata da parcela em contexto web.
Os 36 estudos com dados extraídos foram publicados entre 2010 e 2026, com ao menos um estudo em 16 dos 17 anos do intervalo; só 2017 não tem nenhum. O ano vem do campo ``Ano de Publicação'' da planilha de extração. A Figura~\ref{fig:publicacoes-ano-web} mostra a distribuição.

\begin{figure}[H]
    \centering
    \caption{Publicações por ano e contexto web}
    \label{fig:publicacoes-ano-web}
    \fbox{\includegraphics[width=0.8\textwidth]{figuras/publicacoes-ano-web.png}}
    \vspace{0.2cm}
    \fonte{Autor (2026).}
\end{figure}

Em síntese, os estudos se distribuem de forma estável ao longo de 17 anos, com média de 2,1 por ano e sem crescimento contínuo, e 17 deles (47\%) são de 2020 em diante. Dos 36, só 3 (8\%) têm relação com a web, e apenas um, de 2021, avalia suítes de testes de aplicações web. Isso tem duas implicações para este trabalho. A primeira é que as respostas às demais questões de pesquisa se apoiam quase só em evidência de outros contextos e precisam ser verificadas em aplicações web. A segunda é que o tema segue ativo, mas sem estudos recentes sobre a web no conjunto, o que reforça a relevância de aplicar e avaliar o teste de mutação em aplicações web de código aberto na segunda etapa, a ser detalhada no Capítulo~\ref{cap:planejamento-tcc2}.

\subsection{Em quais revistas e conferências são publicados estudos sobre o tema?}
\label{subsec:revistas-conferencias}

Os 36 estudos com dados extraídos estão em 22 fontes: 6 periódicos, 13 conferências e 3 \textit{workshops}. A Tabela~\ref{tab:fontes-publicacao} reúne as fontes por tipo.

\begin{table}[H]
    \centering
    \scriptsize
    \caption{Fontes de publicação dos estudos}
    \label{tab:fontes-publicacao}
    \begin{tabular}{|p{2.3cm}|p{5.5cm}|c|p{4.7cm}|}
        \hline
        \textbf{Tipo} & \textbf{Fonte} & \textbf{Estudos} & \textbf{Referências} \\
        \hline
        Periódico & IEEE Transactions on Software Engineering & 3 & \cite{zhang2019, zhang2022, guizzo2022} \\
        \hline
        Periódico & Journal of Electronic Testing & 2 & \cite{arasteh2022, arasteh2023} \\
        \hline
        Periódico & ACM Transactions on Software Engineering and Methodology & 1 & \cite{khanfir2023} \\
        \hline
        Periódico & Journal of Systems and Software & 1 & \cite{tao2025} \\
        \hline
        Periódico & IEEE Access & 1 & \cite{lin2018} \\
        \hline
        Periódico & Data Technologies and Applications & 1 & \cite{shomali2020} \\
        \hline
        \textbf{Subtotal} & \textbf{Periódicos} & \textbf{9 (25\%)} & -- \\
        \hline
        Conferência & ESEC/FSE (Foundations of Software Engineering) & 5 & \cite{inozemtseva2013, just2014, kurtz2016, ahmed2016, ojdanic2023} \\
        \hline
        Conferência & ICSE (International Conference on Software Engineering) & 4 & \cite{zhang2010, yao2014, tufano2020, petrovic2021} \\
        \hline
        Conferência & ICST (International Conference on Software Testing, Verification and Validation) & 2 & \cite{nica2011, just2012} \\
        \hline
        Conferência & QRS (International Conference on Software Quality, Reliability and Security) & 2 & \cite{ashong2025, ashong2026} \\
        \hline
        Conferência & SBQS (Simpósio Brasileiro de Qualidade de Software) & 2 & \cite{kuroishi2020, guerino2024} \\
        \hline
        Conferência & ISSTA (International Symposium on Software Testing and Analysis) & 1 & \cite{grano2019} \\
        \hline
        Conferência & ICSME (International Conference on Software Maintenance and Evolution) & 1 & \cite{yandrapally2021} \\
        \hline
        Conferência & APSEC (Asia Pacific Software Engineering Conference) & 1 & \cite{kintis2010} \\
        \hline
        Conferência & CEE-SECR (Central and Eastern European Software Engineering Conference) & 1 & \cite{dimov2010} \\
        \hline
        Conferência & ICACEA (International Conference on Advances in Computer Engineering and Applications) & 1 & \cite{chauhan2015} \\
        \hline
        Conferência & UKSim (International Conference on Computer Modelling and Simulation) & 1 & \cite{assylbekov2013} \\
        \hline
        Conferência & OSSCOM (International Conference on Open Source Software Computing) & 1 & \cite{alenezi2016} \\
        \hline
        Conferência & SDPC (International Conference on Sensing, Diagnostics, Prognostics, and Control) & 1 & \cite{he2024} \\
        \hline
        \textbf{Subtotal} & \textbf{Conferências} & \textbf{23 (64\%)} & -- \\
        \hline
        Workshop & ICSTW (ICST Workshops) & 2 & \cite{esterobotaro2010, ashong2026} \\
        \hline
        Workshop & COMPSAC Workshops (Computer Software and Applications Conference) & 1 & \cite{belli2011} \\
        \hline
        Workshop & RAISE (International Workshop on Realizing AI Synergies in Software Engineering) & 1 & \cite{jalbert2012} \\
        \hline
        \textbf{Subtotal} & \textbf{Workshops} & \textbf{4 (11\%)} & -- \\
        \hline
        \textbf{Total} & \textbf{22 fontes} & \textbf{36 (100\%)} & -- \\
        \hline
    \end{tabular}
    \legend{Fonte: Autor(2026).}
\end{table}

Como é possível perceber, os 36 estudos estão em 22 fontes, com 64\% em conferências, 25\% em periódicos e 11\% em \textit{workshops}. Quatro fontes, o ESEC/FSE, o ICSE, a série ICST e o \textit{IEEE Transactions on Software Engineering}, reúnem 16 estudos (44\%), e as outras 17 têm até dois estudos cada. O uso de periódicos cresce depois de 2018, mas as conferências seguem como o principal meio. Para este trabalho, esses são os eventos e periódicos em que vale buscar estudos complementares e divulgar os resultados. Como a distribuição depende da \textit{string} de busca e dos critérios de seleção da revisão (Seção~\ref{sec:realizacao-busca}), ela não representa todas as fontes que publicam sobre o tema.

\subsection{Quais métodos de validação empírica são utilizados nos estudos?}
\label{subsec:metodos-validacao}

Os 37 estudos foram classificados em sete grupos, segundo o método principal de validação empírica (Tabela~\ref{tab:validacao-empirica}), e um mesmo estudo pode somar técnicas complementares, como testes estatísticos e inspeção manual (Tabela~\ref{tab:tecnicas-complementares}). Em oito estudos o campo ``Métodos de validação empírica'' estava vazio ou trazia só uma expressão, como ``Estudo de Caso'' ou ``Experimento'' \cite{just2012, kintis2010, tufano2020, ahmed2016, alenezi2016, zhang2022, ashong2025, ashong2026}, e o método foi completado a partir dos artigos quando eles estavam disponíveis. A maior parte dos estudos não é de aplicações web, e por isso a análise considera todos eles, e a Seção~\ref{subsubsec:validacao-web} trata dos três em contexto web.

\begin{table}[H]
    \centering
    \scriptsize
    \caption{Métodos de validação empírica, por método principal}
    \label{tab:validacao-empirica}
    \begin{tabular}{|p{3.0cm}|p{3.8cm}|c|c|p{4.0cm}|}
        \hline
        \textbf{Método} & \textbf{Descrição} & \textbf{Estudos} & \textbf{\%} & \textbf{Referências} \\
        \hline
        Experimento comparativo com \textit{benchmarks} ou programas pequenos & Técnica ou estratégia comparada com referência em programas conhecidos ou de pequeno porte & 10 & 28\% & \cite{shomali2020, arasteh2022, arasteh2023, kurtz2016, zhang2010, kuroishi2020, kintis2010, lin2018, tao2025, nica2011} \\
        \hline
        Estudo empírico em projetos reais de código aberto & Vários projetos ou versões, em comparação de técnicas ou em análise observacional & 9 & 25\% & \cite{zhang2019, zhang2022, ashong2025, ashong2026, guizzo2022, ahmed2016, ojdanic2023, alenezi2016} \\
        \hline
        Estudo de caso ou estudo piloto & Aplicação a um ou poucos sistemas, declarada como estudo de caso ou piloto ou restrita a um sistema & 8 & 22\% & \cite{just2012, jalbert2012, belli2011, esterobotaro2010, assylbekov2013, inozemtseva2013, chauhan2015, he2024} \\
        \hline
        Métodos mistos com análise manual & Experimento combinado com inspeção manual, análise de relatos de \textit{bugs} ou questionário & 3 & 8\% & \cite{yao2014, yandrapally2021, guerino2024} \\
        \hline
        Demonstração ou avaliação preliminar & Ferramenta, plataforma ou proposta apresentada com avaliação parcial ou planejada & 3 & 8\% & \cite{tufano2020, dimov2010, grano2019} \\
        \hline
        Validação com falhas reais & Mutantes ou falhas injetadas comparados a falhas reais & 2 & 6\% & \cite{just2014, khanfir2023} \\
        \hline
        Estudo industrial longitudinal & Observação ao longo de anos em ambiente industrial, com grupo de controle & 1 & 3\% & \cite{petrovic2021} \\
        \hline
        \textbf{Total} & \textbf{--} & \textbf{36} & \textbf{100\%} & \textbf{--} \\
        \hline
    \end{tabular}
    \legend{Fonte: Autor (2026).}
\end{table}

\subsubsection{Experimentos comparativos com \textit{benchmarks} e programas pequenos}
\label{subsubsec:experimentos-comparativos}

Dez estudos (28\%) comparam técnicas ou estratégias em \textit{benchmarks} ou programas de pequeno porte. \citeonline{shomali2020} comparam a mutação de todo o código com a de locais selecionados em cinco programas, e \citeonline{arasteh2022} e \citeonline{arasteh2023} usam cinco programas cada, com 10 execuções independentes por algoritmo e comparação com algoritmos concorrentes; no primeiro, os programas têm de 17 a 72 linhas. \citeonline{zhang2010} comparam cinco técnicas de seleção de mutantes em sete programas C, com 50 suítes por configuração, e \citeonline{kurtz2016} fazem um experimento \textit{in vitro} no \textit{benchmark} Siemens, com 512 testes e simulação de Monte Carlo. \citeonline{kuroishi2020} reutilizam programas C e suítes de estudos anteriores e comparam a seleção com a mutação completa e com a amostragem aleatória. \citeonline{kintis2010} examinam estratégias de segunda ordem e a mutação fraca em quinze programas, \citeonline{lin2018} analisam 22 programas e cerca de 96 mil mutantes gerados pelo JCMuta, e \citeonline{tao2025} combinam \textit{benchmarks} e programas descritos como industriais. \citeonline{nica2011} valida a abordagem de casos de teste distinguidores em programas Java de até 100 linhas. Esse método permite controlar as condições e repetir o experimento, mas os programas são pequenos ou de uso conhecido, o que limita a generalização.

\subsubsection{Estudos em projetos reais de código aberto}
\label{subsubsec:estudos-projetos-reais}

Nove estudos (25\%) avaliam projetos reais de código aberto, em geral vários. \citeonline{zhang2019} usam nove projetos-base e 154 projetos adicionais do GitHub, de 172 a 92.176 linhas, com o modelo treinado em uma versão e avaliado na seguinte. \citeonline{zhang2022} usam 163 conjuntos de dados do PMT original e 650 do PMT estendido, todos gerados com o PIT. \citeonline{ashong2025} avaliam 15 projetos Java, de 442 a 48.115 linhas, e \citeonline{ashong2026}, 14 projetos Java, com 1.000 repetições por técnica estocástica e taxa de amostragem. \citeonline{guizzo2022} comparam estratégias em dez sistemas e quatro versões, com 30 execuções independentes por combinação. \citeonline{ahmed2016} relacionam cobertura e \textit{mutation score} às correções de \textit{bugs} posteriores em um grande número de programas Java do GitHub e do Apache, e \citeonline{ojdanic2023} fazem uma análise longitudinal do histórico de projetos de código aberto. \citeonline{alenezi2016} conduzem experimentos em seis projetos Java, e \citeonline{ashong2026} descrevem um experimento, sem mais detalhes na planilha. Esses estudos têm maior validade externa que os experimentos com programas pequenos, mas o custo da análise de mutação limita o número de projetos (Seção~\ref{subsubsec:custo-computacional}).

\subsubsection{Estudos de caso, estudos piloto e demonstrações}
\label{subsubsec:estudos-caso-piloto}

Oito estudos (22\%) são estudos de caso ou estudos piloto. \citeonline{just2012} fazem um estudo de caso com quatro programas reais, de 3.000 a quase 40.000 linhas. \citeonline{jalbert2012} fazem uma avaliação preliminar, que chamam de estudo de caso, no JGAP 3.6.1, com validação cruzada de dez partições. \citeonline{belli2011} usam dois módulos de uma aplicação comercial web, \citeonline{esterobotaro2010} usam três composições WS-BPEL de 110, 202 e 633 linhas, com formalização teórica de apoio, e \citeonline{assylbekov2013} usam um programa, o Spojo, removendo os casos de teste um a um. \citeonline{inozemtseva2013} fazem um estudo piloto com um projeto real, \citeonline{chauhan2015} usam um sistema de informação de estudantes com 10 módulos, e \citeonline{he2024} comparam mutantes de três ordens em um software embarcado, com os mesmos cinco dados de teste. Três estudos (8\%) são demonstrações ou avaliações preliminares. \citeonline{tufano2020} executam o DeepMutation em quatro projetos do Defects4J, como avaliação preliminar da ferramenta. \citeonline{dimov2010} aplicam sua plataforma aos subsistemas de busca e indexação do Apache Lucene, sem consolidação experimental com repetições, intervalos de confiança ou comparação estatística. \citeonline{grano2019} tem parte do estudo realizada, a da pergunta sobre fatores de qualidade dos testes, e parte planejada, com modelos de predição e uma implementação no EvoSuite avaliada em classes do SF110.

\subsubsection{Falhas reais, estudo industrial e métodos mistos}
\label{subsubsec:falhas-reais-industrial}

Poucos estudos usam falhas reais ou ambiente industrial como referência. \citeonline{just2014} respondem a três questões de pesquisa: se falhas reais são acopladas a mutantes ($\chi^2$), que tipos de falha não são acopladas (revisão manual qualitativa de cada falha) e se a detecção de mutantes se correlaciona com a de falhas reais (correlação ponto-bisserial e \textit{rank-bisserial}, Wilcoxon e tamanho de efeito $\hat{A}_{12}$). \citeonline{khanfir2023} comparam o iBiR e a mutação tradicional com amostragem aleatória em falhas reais do Defects4J, e o iBiR distinguiu estatisticamente as suítes que detectam a falha real das que não detectam em 126 falhas, contra 55 da abordagem comparada. \citeonline{petrovic2021} fazem o único estudo industrial, com desenho longitudinal e interventivo ao longo de seis anos, mais de 14 milhões de mutantes e grupos de tratamento e de controle, analisados no nível de arquivos, com correlação de Spearman e teste de Wilcoxon. O \textit{pipeline} de acoplamento partiu de 1.765 \textit{bugs} candidatos, reteve 1.502, gerou cerca de 400 mil mutantes e passou de 33 milhões de execuções, e incluiu a inspeção manual de 50 \textit{bugs} não acoplados. Três estudos combinam métodos: \citeonline{yao2014} usam 18 programas C, inspeção manual sistemática de 1.230 mutantes não mortos por suítes com 100\% de cobertura de ramos factíveis e correlação de Spearman; \citeonline{guerino2024} combinam uma comparação estática, documental e por questionário, com uma comparação dinâmica, em programas escolhidos por conveniência, com geração e ampliação de testes e análise manual de equivalentes; e \citeonline{yandrapally2021}, descritos na Seção~\ref{subsubsec:validacao-web}.

\subsubsection{Técnicas complementares}
\label{subsubsec:tecnicas-complementares}

Além do método principal, 18 estudos (50\%) informam ao menos uma técnica complementar, reunidas na Tabela~\ref{tab:tecnicas-complementares}. A mais frequente é a análise estatística, em dez estudos, com correlação de Spearman em três \cite{yao2014, kurtz2016, petrovic2021}, teste de Wilcoxon em três \cite{just2014, grano2019, petrovic2021} e medidas de tamanho de efeito em dois \cite{just2014, grano2019}. \citeonline{arasteh2023} usam ANOVA, e \citeonline{ashong2026} usam a contagem de Borda e o $\tau$ de Kendall. A inspeção manual aparece em sete estudos, principalmente para classificar mutantes equivalentes, e a repetição de execuções ou de construções de suítes, em seis. Como las técnicas só foram contadas quando a planilha ou o artigo as descrevem, os números podem estar subestimados.

\begin{table}[H]
    \centering
    \scriptsize
    \caption{Técnicas complementares de validação (um estudo pode aparecer em mais de uma)}
    \label{tab:tecnicas-complementares}
    \begin{tabular}{|p{3.2cm}|c|c|p{7.5cm}|}
        \hline
        \textbf{Técnica} & \textbf{Estudos} & \textbf{\%} & \textbf{Referências} \\
        \hline
        Análise estatística & 10 & 28\% & \cite{yao2014, petrovic2021, kurtz2016, just2014, grano2019, arasteh2023, ashong2026, guizzo2022, ahmed2016, khanfir2023} \\
        \hline
        Inspeção manual & 7 & 19\% & \cite{yao2014, yandrapally2021, petrovic2021, just2014, guerino2024, esterobotaro2010, kuroishi2020} \\
        \hline
        Repetição de execuções ou de construções de suítes & 6 & 17\% & \cite{zhang2010, kurtz2016, arasteh2022, arasteh2023, guizzo2022, ashong2026} \\
        \hline
        Falhas reais como referência & 3 & 8\% & \cite{just2014, khanfir2023, petrovic2021} \\
        \hline
        Validação cruzada ou treino e teste em versões distintas & 2 & 6\% & \cite{jalbert2012, zhang2019} \\
        \hline
        Disponibilização pública de dados ou ferramentas & 2 & 6\% & \cite{yandrapally2021, ashong2026} \\
        \hline
    \end{tabular}
    \legend{Fonte: Autor (2026).}
\end{table}

\subsubsection{Métodos nos estudos em contexto web}
\label{subsubsec:validacao-web}

Os três estudos em contexto web (Seção~\ref{subsubsec:publicacoes-web}) usam estudos de caso ou métodos mistos, com poucos objetos. \citeonline{yandrapally2021} analisam manualmente 250 relatos de \textit{bugs}, de 6.331 coletados em 10 aplicações de código aberto, para derivar as categorias de falhas e os operadores, e cerca de 44\% deles afetam a interface. Em seguida, fazem um experimento em seis aplicações (\textit{AddressBook}, \textit{Claroline}, \textit{Collabtive}, \textit{MantisBT}, \textit{MRBS} e \textit{PPMA}), cada uma com suíte JUnit e Selenium, com 50 mutantes por aplicação, 300 no total, de cerca de 586 mil candidatos. Cada mutante é analisado manualmente quanto à equivalência, à severidade e à teimosia, para responder a duas questões: a eficiência em gerar mutantes não equivalentes (93\% em média) e a utilidade dos mutantes para melhorar as suítes (\textit{mutation score} médio de 0,20). \citeonline{belli2011} fazem um estudo de caso com comparação experimental e análise estatística em dois módulos de uma aplicação comercial web, de um sistema de cerca de 53 mil linhas. \citeonline{esterobotaro2010} aplicam um procedimento a três composições WS-BPEL: geram todos os mutantes, executam contra uma suíte inicial feita à mão, inspecionam os sobreviventes, completam a suíte até torná-la adequada e a reduzem a um conjunto não redundante. Nenhum deles usa falhas reais como referência nem muitos projetos, ao contrário dos nove estudos em projetos reais de código aberto (Seção~\ref{subsubsec:estudos-projetos-reais}), e os três têm seis aplicações, dois módulos e três composições como objetos de estudo.

\subsubsection{Síntese e implicações}
\label{subsubsec:sintese-validacao}

Em síntese, os métodos mais usados são o experimento comparativo com \textit{benchmarks} ou programas pequenos (28\%), o estudo em projetos reais de código aberto (25\%) e o estudo de caso ou piloto (22\%). Apenas três estudos (8\%) usam falhas reais ou ambiente industrial como referência, e metade dos estudos informa alguma técnica complementar, em geral estatística ou inspeção manual. Para a web, os três estudos têm poucos objetos. Para este trabalho, isso tem duas implicações. A primeira é que a mineração de repositórios, com seis projetos web de código aberto, se aproxima do segundo grupo, e seis projetos é um número próximo ao de \citeonline{alenezi2016} e de \citeonline{yandrapally2021} e menor que o de \citeonline{guizzo2022}, com dez, e o de \citeonline{ashong2025, ashong2026}, com 14 e 15. A segunda é que a classificação manual dos mutantes sobreviventes, prevista para a segunda etapa, tem precedentes em \citeonline{yao2014}, \citeonline{yandrapally2021} e \citeonline{guerino2024}, que usam análise manual de equivalência, e convém reportar também a concordância entre avaliadores, decisão a ser confirmada no capítulo de proposta.

\subsection{Quais são os principais autores e grupos de pesquisa que investigam o tema?}
\label{subsec:autores-grupos}

Os 37 estudos somam 147 autorias de 123 autores distintos, com média de 4,1 autores por estudo, de um a sete; dois estudos têm um único autor \cite{nica2011, grano2019}. Os autores foram contados pelos nomes do campo ``Autores'' da planilha, sem verificar homônimos, e os grupos foram definidos pela coautoria entre os estudos, com as afiliações tiradas dos próprios artigos quando eles estavam disponíveis. Dezessete autores (14\%) aparecem em mais de um estudo, e 106 (86\%), em apenas um (Tabela~\ref{tab:autores-frequentes}). Os grupos são apresentados na Tabela~\ref{tab:grupos-pesquisa}, e a Seção~\ref{subsubsec:autores-web} trata do contexto web.

\begin{table}[H]
    \centering
    \scriptsize
    \caption{Autores com mais de um estudo}
    \label{tab:autores-frequentes}
    \begin{tabular}{|p{3.2cm}|c|p{7.5cm}|}
        \hline
        \textbf{Autor} & \textbf{Estudos} & \textbf{Referências} \\
        \hline
        Just, René & 3 & \cite{just2012, just2014, petrovic2021} \\
        \hline
        Harman, Mark & 3 & \cite{yao2014, zhang2019, ojdanic2023} \\
        \hline
        Papadakis, Mike & 3 & \cite{kintis2010, ojdanic2023, khanfir2023} \\
        \hline
        Arasteh, Bahman & 3 & \cite{shomali2020, arasteh2022, arasteh2023} \\
        \hline
        Ashong, Magdalene; Laurent, Thomas; Ventresque, Anthony & 3 cada & \cite{ashong2025, ashong2026} \\
        \hline
        Jia, Yue & 2 & \cite{yao2014, zhang2019} \\
        \hline
        Zhang, Lu & 2 & \cite{zhang2010, zhang2019} \\
        \hline
        Yao, Xiangjuan & 2 & \cite{yao2014, tao2025} \\
        \hline
        Fraser, Gordon & 2 & \cite{just2014, petrovic2021} \\
        \hline
        Inozemtseva, Laura; Holmes, Reid & 2 cada & \cite{inozemtseva2013, just2014} \\
        \hline
        Delamaro, Márcio Eduardo & 2 & \cite{kurtz2016, kuroishi2020} \\
        \hline
        Vincenzi, Auri Marcelo Rizzo; Kuroishi, Pedro Henrique & 2 cada & \cite{kuroishi2020, guerino2024} \\
        \hline
        Gharehchopogh, Farhad Soleimanian & 2 & \cite{arasteh2022, arasteh2023} \\
        \hline
    \end{tabular}
    \legend{Fonte: Autor (2026).}
\end{table}

\begin{table}[H]
    \centering
    \scriptsize
    \caption{Grupos de pesquisa, definidos pela coautoria entre os estudos}
    \label{tab:grupos-pesquisa}
    \begin{tabular}{|p{2.5cm}|p{2.6cm}|p{3.2cm}|c|p{3.2cm}|}
        \hline
        \textbf{Grupo} & \textbf{Autores principais} & \textbf{Afiliações nos artigos} & \textbf{Estudos} & \textbf{Referências} \\
        \hline
        Harman, Papadakis e Jia & Harman, Papadakis, Jia, Yao, L. Zhang & University College London (CREST); Athens University of Economics and Business; SnT, University of Luxembourg; Peking University & 7 (19\%) & \cite{yao2014, zhang2019, ojdanic2023, kintis2010, khanfir2023, zhang2010, tao2025} \\
        \hline
        Just, Fraser e Holmes & Just, Fraser, Holmes, Inozemtseva, Kapfhammer & Ulm University e Allegheny College (2012); University of Washington, University of Waterloo e University of Sheffield (2014); Google, University of Passau e University of Washington (2021) & 4 (11\%) & \cite{just2012, inozemtseva2013, just2014, petrovic2021} \\
        \hline
        Ashong, Laurent e Ventresque & Ashong, Laurent, Ventresque & School of Computer Science and Statistics, Trinity College Dublin, e Lero Research Centre & 3 (8\%) & \cite{ashong2025, ashong2026} \\
        \hline
        Arasteh e Gharehchopogh & Arasteh, Gharehchopogh & Istinye University; Islamic Azad University (Urmia Branch) & 3 (8\%) & \cite{shomali2020, arasteh2022, arasteh2023} \\
        \hline
        Delamaro, Vincenzi e Kuroishi & Delamaro, Vincenzi, Kuroishi & Universidade de São Paulo (Delamaro, em \citeonline{kurtz2016}); demais afiliações não verificadas & 3 (8\%) & \cite{kurtz2016, kuroishi2020, guerino2024} \\
        \hline
        Sem coautoria com outros estudos & -- & -- & 16 (44\%) & \cite{jalbert2012, chauhan2015, nica2011, esterobotaro2010, belli2011, lin2018, zhang2022, dimov2010, grano2019, yandrapally2021, he2024, tufano2020, alenezi2016, assylbekov2013, ahmed2016, guizzo2022} \\
        \hline
        \textbf{Total} & \textbf{--} & \textbf{--} & \textbf{36 (100\%)} & \textbf{--} \\
        \hline
    \end{tabular}
    \legend{Fonte: Autor(2026).}
\end{table}

\subsubsection{Autores mais frequentes}
\label{subsubsec:autores-frequentes}

Sete autores aparecem em três estudos: Just, Harman, Papadakis e Arasteh, e os três autores que assinam juntos os estudos de \citeonline{ashong2025, ashong2025b, ashong2026}. Dez autores aparecem em dois. Os autores recorrentes publicam ao longo de vários anos: Papadakis aparece de 2010 a 2023 \cite{kintis2010, khanfir2023}, Just de 2012 a 2021 \cite{just2012, petrovic2021}, Harman de 2014 a 2023 \cite{yao2014, ojdanic2023} e Arasteh de 2020 a 2023 \cite{shomali2020, arasteh2023}. Harman e Papadakis assinam juntos um estudo, o de \citeonline{ojdanic2023}. Lu Zhang foi considerado o mesmo autor nos dois estudos em que aparece \cite{zhang2010, zhang2019}, pelo nome.

\subsubsection{Grupos de pesquisa}
\label{subsubsec:grupos-pesquisa}

Cinco grupos reúnem 20 dos 37 estudos (56\%), e os outros 16 (44\%) são de equipes que não se repetem no conjunto. A maior rede é a de Harman, Papadakis e Jia, com sete estudos (19\%), ligados por coautoria: \citeonline{yao2014} tratam de mutantes equivalentes e teimosos, \citeonline{zhang2019} da predição do \textit{mutation score}, \citeonline{kintis2010} de alternativas ao teste de mutação, \citeonline{ojdanic2023} de conjuntos de mutantes ao longo da evolução do código, \citeonline{khanfir2023} da injeção de falhas guiada por relatos de \textit{bugs}, \citeonline{zhang2010} da seleção de mutantes e \citeonline{tao2025}, por meio de Yao, da geração de dados de teste. As afiliações nos artigos são a University College London (CREST), a Athens University of Economics and Business, o SnT da University of Luxembourg e a Peking University. O grupo de Just, Fraser e Holmes tem quatro estudos (11\%): mutantes redundantes \cite{just2012}, seleção por histórico de defeitos \cite{inozemtseva2013}, mutantes como substitutos de falhas reais \cite{just2014} e adoção no Google \cite{petrovic2021}. As afiliações mudam entre os artigos, de Ulm e Allegheny College para Washington, Waterloo e Sheffield e depois para Google e Passau, e o grupo se mantém pela coautoria, e não pela instituição. Os outros três grupos têm três estudos cada (8\%). Ashong, Laurent e Ventresque, do Trinity College Dublin e do Lero Research Centre, tratam da seleção estática de mutantes em 2025 e 2026. Arasteh e Gharehchopogh, da Istinye University e da Islamic Azad University, reduzem mutantes com algoritmos metaheurísticos \cite{shomali2020, arasteh2022, arasteh2023}. Delamaro, Vincenzi e Kuroishi têm estudos sobre mutantes dominadores, com Ammann e Offutt, da George Mason University \cite{kurtz2016}, sobre arcos primitivos \cite{kuroishi2020} e sobre ferramentas para Python \cite{guerino2024}; os dois estudos do SBQS (Seção~\ref{subsubsec:conferencias-workshops}) são deste grupo. Entre as equipes que não se repetem estão \citeonline{yandrapally2021} (University of British Columbia), \citeonline{tufano2020} (Washington and Lee University, Università della Svizzera italiana e University of Sannio), \citeonline{ahmed2016} (Oregon State University), \citeonline{grano2019} (University of Zurich), \citeonline{lin2018} (Beijing University of Posts and Telecommunications), \citeonline{esterobotaro2010} (University of Cádiz), \citeonline{nica2011} (Graz University of Technology), \citeonline{assylbekov2013} (University of Melbourne) e \citeonline{alenezi2016} (Prince Sultan University e Yarmouk University).

\subsubsection{Autores e grupos nos estudos em contexto web}
\label{subsubsec:autores-web}

Os três estudos em contexto web (Seção~\ref{subsubsec:publicacoes-web}) são de equipes diferentes e nenhum deles tem autor entre os 17 que se repetem. \citeonline{yandrapally2021} são da University of British Columbia, \citeonline{esterobotaro2010} são da University of Cádiz, e as afiliações de \citeonline{belli2011} não foram verificadas. Nenhum dos cinco grupos tem estudo sobre aplicações web no conjunto, e nenhum grupo tem mais de um estudo em contexto web. Assim, o tema web aparece em estudos isolados, e não em uma linha de pesquisa continuada pelos mesmos autores no conjunto de estudos.

\subsubsection{Síntese e implicações}
\label{subsubsec:sintese-autores}

Em síntese, os 37 estudos têm 123 autores, dos quais 17 (14\%) aparecem em mais de um estudo, e cinco grupos reúnem 20 estudos (56\%), com a rede de Harman, Papadakis e Jia como a maior, com sete. A concentração é moderada, pois quase metade dos estudos (44\%) é de equipes que não se repetem. Nenhum grupo se dedica à web no conjunto: os três estudos em contexto web são de autores isolados. Para este trabalho, os grupos recorrentes são referências para fundamentar e comparar os resultados e para buscar outras publicações, e a escassez de grupos na web reforça a relevância de estudar mutantes em aplicações web. Como os grupos foram definidos só pela coautoria entre os estudos da revisão, podem existir vínculos fora do conjunto, e as filiações foram verificadas apenas nos artigos disponíveis.